% concatenated from final-version-0109.tex



\documentclass[11pt,a4paper,twoside]{amsart}
\usepackage{amsmath}
\usepackage{amssymb}
\usepackage{amsfonts}
\usepackage{amscd}
\usepackage{xcolor}
\usepackage{url}
\newtheorem{theorem}{Theorem}[section]
\newtheorem{lemma}[theorem]{Lemma}
\newtheorem{proposition}[theorem]{Proposition}
\newtheorem{corollary}[theorem]{Corollary}
\theoremstyle{definition}
\newtheorem{definition}[theorem]{Definition}
\newtheorem{claim}[theorem]{Claim}
\newtheorem{example}[theorem]{Example}
\newtheorem{question}[theorem]{Question}
\newtheorem{questions}[theorem]{Questions}
\newtheorem{conjecture}[theorem]{Conjecture}
\theoremstyle{remark}
\newtheorem{remark}[theorem]{Remark}
\newtheorem{note}[theorem]{Note}

\newcommand{\tori}{\mathfrak t}
\newcommand{\mB}{\mathbb B}
\newcommand{\mC}{\mathbb C}
\newcommand{\mD}{\mathbb D}
\newcommand{\mE}{\mathbb E}
\newcommand{\mF}{\mathbb F}
\newcommand{\mG}{\mathbb G}
\newcommand{\mH}{\mathbb H}
\newcommand{\mI}{\mathbb I}
\newcommand{\mJ}{\mathbb J}
\newcommand{\mK}{\mathbb K}
\newcommand{\mL}{\mathbb L}
\newcommand{\mM}{\mathbb M}
\newcommand{\mN}{\mathbb N}
\newcommand{\mO}{\mathbb O}
\newcommand{\mP}{\mathbb P}
\newcommand{\mQ}{\mathbb Q}
\newcommand{\mR}{\mathbb R}
\newcommand{\mS}{\mathbb S}
\newcommand{\mT}{\mathbb T}
\newcommand{\mU}{\mathbb U}
\newcommand{\mV}{\mathbb V}
\newcommand{\mW}{\mathbb W}
\newcommand{\mX}{\mathbb X}
\newcommand{\mY}{\mathbb Y}
\newcommand{\mA}{\mathbb A}
\newcommand{\mZ}{\mathbb Z}
\newcommand{\bK}{\mathbf K}
\newcommand{\bA}{\mathbf A}
\newcommand{\bX}{\mathbf X}
\newcommand{\bU}{\mathbf U}
\newcommand{\bW}{\mathbf W}
\newcommand{\bV}{\mathbf V}
\newcommand{\bZ}{\mathbf Z}
\newcommand{\bY}{\mathbf Y}
\newcommand{\bF}{\mathbf F}
\newcommand{\bP}{\mathbf P}
\newcommand{\ho}{\hookrightarrow}
\newcommand{\Gg}{\gamma}
\newcommand{\GG}{\Gamma}
\newcommand{\bO}{\Omega}
\newcommand{\bo}{\omega}
\newcommand{\bl}{\lambda}
\newcommand{\bL}{\Lambda}
\newcommand{\bs}{\sigma}
\newcommand{\bS}{\Sigma}
\newcommand{\ep}{\epsilon}
\newcommand{\D}{\Delta}
\newcommand{\z}{\zeta}
\newcommand{\kk}{\kappa}
\newcommand{\bsl}{\backslash}
\newcommand{\mcA}{\mathcal A}
\newcommand{\mcB}{\mathcal B}
\newcommand{\mcC}{\mathcal C}
\newcommand{\mcD}{\mathcal D}
\newcommand{\mcE}{\mathcal E}
\newcommand{\mcF}{\mathcal F}
\newcommand{\mcG}{\mathcal G}
\newcommand{\mcH}{\mathcal H}
\newcommand{\mcI}{\mathcal I}
\newcommand{\mcJ}{\mathcal J}
\newcommand{\mcK}{\mathcal K}
\newcommand{\mcL}{\mathcal L}
\newcommand{\mcM}{\mathcal M}
\newcommand{\mcN}{\mathcal N}
\newcommand{\mcO}{\mathcal O}
\newcommand{\mcP}{\mathcal P}
\newcommand{\mcQ}{\mathcal Q}
\newcommand{\mcR}{\mathcal R}
\newcommand{\mcS}{\mathcal S}
\newcommand{\mcT}{\mathcal T}
\newcommand{\mcU}{\mathcal U}
\newcommand{\mcV}{\mathcal V}
\newcommand{\mcW}{\mathcal W}
\newcommand{\mcX}{\mathcal X}
\newcommand{\mcY}{\mathcal Y}
\newcommand{\mcZ}{\mathcal Z}
\newcommand{\CC}{\mathbb C}
\newcommand{\QQ}{\mathbb Q}
\newcommand{\ZZ}{\mathbb Z}
\newcommand{\FF}{\mathbb F}
\newcommand{\RR}{\mathbb R}
\newcommand{\VV}{\mathbb V}
\newcommand{\XX}{\mathbb X}
\newcommand{\YY}{\mathbb Y}
\newcommand{\WW}{\mathbb W}
\newcommand{\PP}{\mathbb P}
\newcommand{\bd}{\bar d}
\newcommand{\fa}{\mathfrak a}
\newcommand{\fb}{\mathfrak b}
\newcommand{\fc}{\mathfrak c}
\newcommand{\fd}{\mathfrak d}
\newcommand{\fe}{\mathfrak e}
\newcommand{\ff}{\mathfrak f}
\newcommand{\fg}{\mathfrak g}
\newcommand{\fh}{\mathfrak h}
\newcommand{\fii}{\mathfrak i}
\newcommand{\fj}{\mathfrak j}
\newcommand{\fk}{\mathfrak k}
\newcommand{\fl}{\mathfrak l}
\newcommand{\fm}{\mathfrak m}
\newcommand{\fn}{\mathfrak n}
\newcommand{\fo}{\mathfrak o}
\newcommand{\fp}{\mathfrak p}
\newcommand{\fq}{\mathfrak q}
\newcommand{\fr}{\mathfrak r}
\newcommand{\fs}{\mathfrak s}
\newcommand{\ft}{\mathfrak t}
\newcommand{\fu}{\mathfrak u}
\newcommand{\fv}{\mathfrak v}
\newcommand{\fw}{\mathfrak w}
\newcommand{\fx}{\mathfrak x}
\newcommand{\fy}{\mathfrak y}
\newcommand{\fz}{\mathfrak z}
\newcommand{\ti}{\widetilde}
\newcommand{\wt}{\widetilde}
\newcommand{\va}{\vartriangleleft}
\newcommand{\un}{\underline}
\newcommand{\ov}{\overline}
\newcommand{\ra}{\rightarrow}
\newcommand{\lr}{\leftrightarrow}
\newcommand{\emp}{\emptyset}
\newcommand{\bFr}{\mathfrak}
\newcommand{\imp}{\Longrightarrow}
\newcommand{\pl}{\varprojlim}
\newcommand{\il}{\varinjlim}
\newcommand{\rt}{\rtimes}
\newenvironment{pf}{\proof[\proofname]}{\endproof}
\newcommand{\supp}{\operatorname{supp}}
\newcommand{\sing}{\operatorname{sing}}
\newcommand{\codim}{\operatorname{codim}}
\newcommand{\sm}{\setminus}
\newcommand{\prl}{\varprojlim}
\newcommand{\inl}{\varinjlim}
\newcommand{\sms}{\smallskip}
\newcommand{\ms}{\medskip}
\newcommand{\Ad}{\operatorname{Ad}}
\definecolor{Red}{rgb}{0.7,0,0}
\newcommand{\red}[1]{{\color{Red}#1}}
\newcommand{\otheta}{\overline{\theta}}
\newcommand{\ind}{\operatorname{ind}}
\newcommand{\Aut}{\operatorname{Aut}}
\newcommand{\St}{\operatorname{St}}
\newcommand{\Sp}{\operatorname{Sp}}
\newcommand{\Id}{\operatorname{Id}}
\newcommand{\Hom}{\operatorname{Hom}}
\newcommand{\End}{\operatorname{End}}
\newcommand{\SL}{\operatorname{SL}}
\newcommand{\opsi}{\overline{\psi}}
\newcommand{\PGL}{\operatorname{PGL}}
\newcommand{\Vect}{\operatorname{Vect}}
\newcommand{\Ind}{\operatorname{Ind}}
\newcommand{\Pro}{\operatorname{Pro}}
\newcommand{\Rep}{\operatorname{Rep}}
\newcommand{\Gr}{\operatorname{Gr}}
\newcommand{\oGr}{\overline{\Gr}}
\newcommand{\oS}{\overline{S}}
\newcommand{\tilR}{\widetilde{R}}
\newcommand{\hZ}{\frac{1}{2}\mathbb Z} 

\newcommand{\bSet}{\mathbf{Set}}

\newcommand{\mSet}{\mathbb{S}\mathrm{et}}
\newcommand{\gr}{\operatorname{gr}}

\newcommand{\mVect}{\mathbb{V}\mathrm{ect}}
\begin{document}
\title[Cuspidal representations over a two-dimensional local field]{An analogue of irreducible cuspidal representations for the group $\PGL(2)$ over a two-dimensional local field}
\author{Alexander Braverman and David Kazhdan}
\date{\today}
\begin{abstract}
Let $F$ be a local non-Archimedean field of odd residue characteristic, let $K=F((t))$, and let $G=\PGL(2)$. We construct analogues of irreducible cuspidal representations of $G(F)$ for the group $G(K)$. The construction starts, up to a metaplectic ambiguity in one even-depth ramified case, from a quadratic extension $L/K$ and a character $\theta:L^*/K^*\to\CC^*$ which is not Galois invariant. We show that the restriction of each representation so obtained to the Borel subgroup $P(K)\subset G(K)$ is irreducible. In contrast with the classical case, these restrictions are not isomorphic to the standard irreducible representation of $P(K)$.
\end{abstract}
\maketitle
\section{Introduction}
\subsection{The setup}
Let $F$ be a local non-Archimedean field of odd residue characteristic with ring of integers $O_F$. Set $K=F((t))$ and $O=F[[t]]$. Let $G$ denote the algebraic group $\PGL(2)$. We denote by $P\subset G$ the Borel subgroup.\footnote{We denote it by $P$ rather than by $B$ because we hope to generalize the results of this paper to $\PGL(n)$ for $n\geq2$, in which case $P$ will be the maximal parabolic subgroup stabilizing a line.} We denote by $T$ the diagonal torus and by $N$ the unipotent radical of $P$.
Let $G'$ be the group scheme over $O$ whose group of $O$-points is the stabilizer, under the adjoint action, of the lattice $Oe\oplus Oh\oplus tOf$, where $e,h,f$ is the standard basis of the Lie algebra $\fg=\mathfrak{sl}_2$. The base change of $G'$ to $K$ is canonically isomorphic to $G$. In particular, $G'(O)$ is canonically a subgroup of $G(K)$, and $G'(O)$ contains $P(O)$.
In this paper we study a class of \emph{special} representations of $G(O)$, $G'(O)$, and $G(K)$. The representation theory of $G(O)$ and $G'(O)$ belongs to the classical representation theory of $p$-adic groups: these groups are projective limits of finite-dimensional algebraic $F$-groups, and we consider representations which factor through smooth representations of finite-dimensional quotients. The notion of a smooth representation of $G(K)$ is subtler. In particular, one must work with actions on pro-vector spaces rather than only on vector spaces; the corresponding category was defined in \cite{GK}. We denote the categories of smooth representations of $G(O)$, $G'(O)$, and $G(K)$ by $\mcR(G(O))$, $\mcR(G'(O))$, and $\mcR(G(K))$. The category $\mcR(G(K))$ has no natural forgetful functor to complex vector spaces; instead, it has a forgetful functor to complex pro-vector spaces.
\begin{definition}\label{spe}
\begin{enumerate}
\item A complex representation $\pi$ of dimension greater than $1$ of $G(O)$ or $G'(O)$ is \emph{special} if its restriction to $P(O)$ is irreducible.
\item A representation $\Pi$ of $G(K)$ is \emph{special} if it is not one-dimensional and its restriction to $P(K)$ is irreducible.
\end{enumerate}
\end{definition}
The main purpose of this paper is to construct special representations of $G(K)$. We conjecture that the construction gives all special representations. To motivate Definition \ref{spe}, we first review the parallel results in the classical situation, in which $K$ is replaced by $F$.
\subsection{Classical results}\label{classical}
All the results of this subsection can be found in \cite{BH}. In this subsection only, $G'$ denotes the group scheme over $O_F$ defined as above, with the uniformizer $\varpi$ in place of $t$.
The embedding $P(F)\ho G(F)$ induces embeddings $P(O_F)\ho G(O_F)$ and $P(O_F)\ho G'(O_F)$. A representation of dimension greater than $1$ of one of the groups $G(O_F)$, $G'(O_F)$, or $G(F)$ is called special if its restriction to $P(O_F)$ or $P(F)$, respectively, is irreducible. An irreducible representation of $G(F)$ is special if and only if it is cuspidal.
\begin{theorem}\label{classical-theorem}
\begin{enumerate}
\item Let $E/F$ be a quadratic extension and let $\theta:E^*/F^*\to\CC^*$ be a character such that $\theta\neq\otheta$, where $\otheta$ is the Galois conjugate of $\theta$. If $E/F$ is unramified, one can attach to $\theta$ a special representation $\pi(E,\theta)$ of $G(O_F)$. If $E/F$ is ramified, one can attach to $(E,\theta)$ a special representation $\pi(E,\theta)$ of $G'(O_F)$.
\item Every special representation of $G(O_F)$ or $G'(O_F)$ is isomorphic to some $\pi(E,\theta)$.
\item Let $\pi$ be an irreducible representation of either $G(O_F)$ or $G'(O_F)$, and let $\Pi(\pi)$ be its compact induction to $G(F)$. Then $\pi$ is special if and only if $\Pi(\pi)$ is irreducible; in this case $\Pi(\pi)$ is special. Thus $\Pi(E,\theta):=\Pi(\pi(E,\theta))$ is irreducible and cuspidal, and every special representation of $G(F)$ is of this form.
\item The representations $\Pi(E_1,\theta_1)$ and $\Pi(E_2,\theta_2)$ are isomorphic if and only if $E_1=E_2$ and either $\theta_1=\theta_2$ or $\theta_1=\otheta_2$. The analogous statement holds for $\pi(E,\theta)$.
\end{enumerate}
\end{theorem}
The third assertion of Theorem \ref{classical-theorem} gives an equivalent characterization of special representations of $G(O_F)$ and $G'(O_F)$: a representation is special precisely when its compact induction to $G(F)$ is irreducible. One could try to use this as the definition after replacing $F$ by $K$, because an analogue of induction is constructed in \cite{GK}. It is more convenient here to use Definition \ref{spe}, since it allows us to work directly with $G(O)$ and $G'(O)$. We do not know whether the two definitions are equivalent for $G(K)$.
\begin{remark}
The representations $\Pi(\pi)$ are unitarizable. Consequently, their irreducibility is equivalent to the equality $\End_{G(F)}(\Pi(\pi))=\CC$.
\end{remark}
\subsection{What is done in this paper?}
We partially extend Theorem \ref{classical-theorem} to $G(O)$, $G'(O)$, and $G(K)$. The paper has two parts. In the first part we study special representations of $G(O)$ and $G'(O)$; in the second part we study their induction to $G(K)$.
\subsection{Special and elliptic representations of $G(O)$ and $G'(O)$}
\subsubsection{Ellipticity}
The definition of a special representation of $G(O)$ or $G'(O)$ is not explicit. Our first result identifies special representations with representations satisfying a more concrete ellipticity condition.\footnote{The notion of ellipticity makes sense for every reductive group.}

Let $G(O)(m)$ be the $m$-th congruence subgroup of $G(O)$, where $m\geq1$, and put $G(O)_m=G(O)/G(O)(m+1)$. For an irreducible representation $\pi$ of $G(O)$, define its depth $m(\pi)$ to be the least $m\geq0$ such that $\pi$ factors through $G(O)_m$.

For $m>0$, the group $G(O)_m$ contains a commutative normal subgroup $G(O)_m(m)$ canonically isomorphic to $\fg(F)$. If $m(\pi)>0$, let $\bo(\pi)\subset\mathfrak{sl}_2(F)$ be the support of the restriction of $\pi$ to this subgroup, using the fixed additive character of $F$ and the trace pairing to identify its character group with $\mathfrak{sl}_2(F)$.
\begin{definition}
An irreducible representation $\pi$ of $G(O)$ is elliptic if either $m(\pi)=0$ and $\pi$ is a cuspidal representation of $G(F)=G(O)_0$, or $m(\pi)>0$ and $\bo(\pi)$ is an elliptic adjoint orbit in $\fg(F)$. The notion of ellipticity for $G'(O)$ is defined similarly; see Definition \ref{ac}.
\end{definition}
We prove that an irreducible representation of $G(O)$ or $G'(O)$ is special if and only if it is elliptic.
\subsubsection{A description of special representations of $G(O)$ and $G'(O)$}
\begin{definition}
A quadratic extension $L/K$ is unramified if $L=E((t))$ for a quadratic extension $E/F$; otherwise it is ramified.
\end{definition}

Let $\widehat G(O)$ and $\widehat G'(O)$ denote the sets of isomorphism classes of special representations of $G(O)$ and $G'(O)$. In the formulation of Theorem \ref{L} we use notation for subsets  $\widehat G(O)_m \subset \widehat G(O) $ and $ \widehat G'(O)_m \subset \widehat G'(O)$  of depth $m$ (see Definition \ref{depth}) and also  for groups $\bar\Xi^L(m)$ and $\ti\Xi^L(m)$ (see Definition \ref{qext1}). 

\begin{theorem}\label{L}
\begin{enumerate}
\item The set $\widehat G(O)$ is the disjoint union of subsets $\widehat G(O)^E$, where $E$ runs through the quadratic extensions of $F$.
\item The set $\widehat G'(O)$ is the disjoint union of subsets $\widehat G'(O)^L$, where $L$ runs through the ramified quadratic extensions of $K$.
\item For every ramified quadratic extension $L/K$, the set $\widehat G'(O)^L$ is naturally in bijection with the characters $\theta:L^*/K^*\to\CC^*$ satisfying $\theta\neq\otheta$, modulo $\theta\sim\otheta$.
\item Let $E/F$ be quadratic and put $L=E((t))$. The depth-zero subset $\widehat G(O)^E_0$ is naturally in bijection with $\bar\Xi^L(0)$.
\item For every positive odd depth $m$, the set $\widehat G(O)^E_m$ is naturally in bijection with $\bar\Xi^L(m)$.
\item If $m>0$ is even and $E/F$ is unramified, then $\widehat G(O)^E_m$ is naturally in bijection with
 $\bar\Xi^L(m)$.
\item If $m>0$ is even and $E/F$ is ramified, then
 $\widehat G(O)^E_m$ is naturally a $\ZZ/2\ZZ$-torsor over $\ti\Xi^L(m)$. 
\end{enumerate}
\end{theorem}
\begin{remark}
In the last case the group  $\D _L $ of quadratic characters of $T^L=L^*/K^*$ is isomorphic to $\ZZ/2\ZZ$ and the map $\bar\Xi^L(m)\to\ti\Xi^L(m)$ is another
 $\ZZ/2\ZZ$-torsor over $\ti\Xi^L(m)$. We do not know whether the  $\ZZ/2\ZZ$-torsors 
 $\widehat G(O)^E_m$ and $ \bar\Xi^L(m) $ over $ \ti\Xi^L(m) $ are canonically isomorphic.
\end{remark}

\subsection{Special representations of $G(K)$}
Our ultimate goal is to construct irreducible representations of $G(K)$. Although we do not have a direct analogue of the usual notion of a cuspidal representation for $G(K)$, Definition \ref{spe} makes sense.
The analogue of induction from $G(O)$ or $G'(O)$, and more generally from a bounded subgroup of $G(K)$, was constructed in \cite{GK}. For a representation $\pi$ of $G(O)$ or $G'(O)$, denote the resulting representation of $G(K)$ by $\Pi(\pi)$.
\begin{theorem}\label{irreducible}
If $\pi$ is a special representation of $G(O)$ or $G'(O)$, then $\Pi(\pi)$ is irreducible. Moreover, its restriction to $P(K)$ is irreducible.
\end{theorem}
Combining Theorems \ref{L} and \ref{irreducible}, we can attach, up to the even-depth ramified ambiguity described above, an irreducible representation $\Pi(L,\theta)$ of $G(K)$ to a quadratic extension $L/K$ and a character $\theta:L^*/K^*\to\CC^*$ satisfying $\theta\neq\otheta$. Replacing $\theta$ by $\otheta$ does not change the resulting representation. We do not prove here that no further identifications occur after induction to $G(K)$; see Section~\ref{quest}.
\subsection{Remarks}
For $G(F)$, the restriction of every special irreducible representation to $P(F)$ is irreducible, and all these restrictions are isomorphic. Indeed, $P(F)$ has a unique irreducible representation of dimension greater than $1$.
This statement fails after replacing $F$ by $K$. One can construct an irreducible representation $\Upsilon$ of $P(K)$ which is a direct analogue of the unique irreducible cuspidal representation of $P(F)$, but $\Upsilon$ does not extend to $G(K)$. We expect that $\Pi(\pi)|_{P(K)}$ depends on $\pi$; more precisely, we expect that representations of different depths have nonisomorphic restrictions. Nevertheless, all representations $\Pi(\pi)$ constructed from special representations of $G(O)$ or $G'(O)$ admit a $\ZZ$-filtration whose associated graded has a natural $P(K)$-action and is isomorphic to $\Upsilon$.
We also expect that every special representation of $G(K)$ is induced from a special representation of $G(O)$ or $G'(O)$, but we do not presently know how to prove this.
Finally, $G(K)$ admits a natural central extension by $F^*$, and one can study representations of this extension at a prescribed level, as in \cite{GK}. The choice of level appears not to affect the special representations considered here, so throughout this paper we work at level zero.
\subsection{A remark about terminology}
We use the word ``Claim'' for statements for which we omit a proof, either because we give a reference or because the verification is straightforward. Statements for which we provide proofs are called theorems, propositions, lemmas, or corollaries.
We also use the phrase ``algebraic $F$-variety $Y$'' for the set of $F$-points $\un Y(F)$ of an algebraic variety $\un Y$ over $F$. Such a set carries its natural topology.
\subsection{Organization of the paper}
Section \ref{deff} contains basic definitions. Section \ref{reptheory} reviews classical representation-theoretic results. Section \ref{PO} classifies smooth irreducible representations of $P(O)$. Sections \ref{GO} and \ref{G'O} classify special representations of $G(O)$ and $G'(O)$. The subsequent part of the paper studies representations of $P(K)$ and the representations of $G(K)$ induced from these special representations.
\subsection{Acknowledgments}
We thank Michael Finkelberg for very helpful discussions. During the preparation of this paper, D.K. was partially supported by ERC grant no.~101142781.
\section{Definitions and notation}\label{deff}
\subsection{Quadratic extensions and tori over general fields}
Let $K$ be a field of characteristic different from $2$, and let $L/K$ be a quadratic extension.
\begin{definition}
\begin{enumerate}
\item Put $T^L=L^*/K^*$. This is naturally the group of $K$-points of an algebraic $K$-torus.
\item Put $L^1=\{l\in L^*:N_{L/K}(l)=1\}$.
\end{enumerate}
\end{definition}
\begin{lemma}\label{j}
\begin{enumerate}
\item Every quadratic extension $L/K$ has the form $L=K(\sqrt a)$ for some $a\in K^*\sm(K^*)^2$.
\item The map $T^L\to L^1$ given by $l\mapsto l/\bar l$ is an isomorphism.
\end{enumerate}
\end{lemma}
\begin{proof}
The first assertion is the usual classification of quadratic extensions in characteristic different from $2$. For the second, the displayed map is well defined on $L^*/K^*$. Its kernel is $K^*$, and Hilbert's Theorem~90 shows that every element of $L^1$ is of the form $l/\bar l$.
\end{proof}
\begin{remark}\label{jj}
We use Lemma \ref{j}(2) to identify $T^L$ with $L^1$.
\end{remark}
\begin{definition}\label{G-K-def}
\begin{enumerate}
\item $G=\PGL(2)$, and $P\subset G$ is the Borel subgroup of upper triangular matrices.
\item For $a\in K^*$, put $\widehat a=\begin{pmatrix}a&0\\0&1\end{pmatrix}\in P(K)$.
\item For $b\in K$, put $\widetilde b=\begin{pmatrix}1&b\\0&1\end{pmatrix}\in P(K)$.
\item The Lie algebra of $G$ is denoted by $\fg$.
\end{enumerate}
\end{definition}
\subsection{Laurent power series}
From now on, assume that $K=F((t))$, where $F$ is a field of characteristic different from $2$.
\begin{definition}\label{intf}
\begin{enumerate}
\item $O=F[[t]]$ and $\fm=tO$.
\item $O_m=O/\fm^{m+1}$ and $O_m(n)=\fm^n/\fm^{m+1}$ for $n\leq m$.
\item $O^*(n)=1+t^nO$ for $n\geq1$.
\item $v:K^*\to\ZZ$ is the valuation normalized by $v(t)=1$.
\end{enumerate}
\end{definition}
\subsubsection{Quadratic extensions}
\begin{definition}
\begin{enumerate}
\item A quadratic extension $K(\sqrt a)/K$ is unramified if $v(a)$ is even and ramified if $v(a)$ is odd.
\item $L^1(1)=\{l\in L^1:v_L(l-1)\geq1\}$, where $v_L$ is the normalized valuation on $L$.
\end{enumerate}
\end{definition}
\begin{claim}\label{quad-LK}
\begin{enumerate}
\item Every unramified quadratic extension $L/K$ has the form $E((t))$, where $E=F(\sqrt a)$ for some $a\in F^*\sm(F^*)^2$. In this case $L^1=E^1\times L^1(1)$.
\item Every ramified quadratic extension $L/K$ has the form $K(\sqrt a)$ with $v(a)=1$, and $L^1=\{\pm1\}\times L^1(1)$.
\end{enumerate}
\end{claim}
\begin{definition}\label{G-O-def}
\begin{enumerate}
\item For $g\in G(K)$, let $v(g)\in\ZZ_{\leq0}$ be the largest integer $r$ such that $\Ad(g)(\fg(O))\subset t^r\fg(O)$.
\item Let $I\subset G(O)$ be the Iwahori subgroup represented by matrices $\begin{pmatrix}a&b\\c&d\end{pmatrix}$ with $c\in\fm$.
\item Put $G'(O)=\ZZ/2\ZZ\ltimes I$, where the nontrivial element is represented by $s=\begin{pmatrix}0&1\\t&0\end{pmatrix}$.
\item A subset $A\subset G(K)$ is bounded if $v(A)$ is bounded below.
\end{enumerate}
\end{definition}
\begin{claim}
The group $G'(O)$ is the normalizer of $I$ in $G(K)$.
\end{claim}
\begin{remark}
We regard $G(O)$ and $G'(O)$ as groups of $F$-points of pro-algebraic groups over $F$.
\end{remark}
\begin{lemma}\label{BT}
\begin{enumerate}
\item Every bounded subgroup of $G(K)$ is conjugate to a subgroup of $G(O)$ or to a subgroup of $G'(O)$.
\item The product maps $P(K)\times G(O)\to G(K)$ and $P(K)\times G'(O)\to G(K)$ are surjective.
\end{enumerate}
\end{lemma}
\begin{proof}
Part (1) follows from Bruhat--Tits theory; see \cite[Section~3.2]{BT}. The first assertion of part (2) is the Iwasawa decomposition. For the second, use the  decomposition
$$
G(K)=P(K)I\sqcup P(K)wI,
$$
where $w$ is a lift of the nontrivial Weyl element. We may take $w=s\in G'(O)$, and $I\subset G'(O)$, whence $G(K)=P(K)G'(O)$.
\end{proof}
\subsection{Tori in $G(O)$ and $G'(O)$}
Let $L/K$ be quadratic and let $O_L$ be its ring of integers. The regular representation of $L^*$ on the two-dimensional $K$-vector space $L$ gives an embedding $L^*\ho\operatorname{GL}_K(L)$ and hence an embedding $T^L\ho G(K)$ after choosing a $K$-basis of $L$.
If $L/K$ is unramified, an $O$-basis of $O_L$ identifies the regular representation of $O_L^*$ with a subgroup of $\operatorname{GL}_2(O)$. If $L/K$ is ramified, choose an $O$-basis adapted to the lattice chain $O_L\supset\mathfrak p_L$; equivalently, identify this chain with the standard chain $O^2\supset Oe_1\oplus tOe_2$. The regular representation then normalizes the standard chain and gives an embedding into $G'(O)$.
\begin{claim}\label{imb}
\begin{enumerate}
\item If $L/K$ is unramified, the resulting subgroup $j^L(T^L)$ lies in $G(O)$, and its $G(O)$-conjugacy class is independent of the chosen $O$-basis of $O_L$.
\item If $L/K$ is ramified, the subgroup obtained from an adapted basis lies in $G'(O)$, and its $G'(O)$-conjugacy class is independent of the adapted basis.
\end{enumerate}
\end{claim}
\begin{proof}
In the unramified case, two $O$-bases differ by an element of $\operatorname{GL}_2(O)$. In the ramified case, two identifications of the lattice chain $O_L\supset\mathfrak p_L$ with the standard lattice chain differ by an element of its normalizer, namely $G'(O)$.
\end{proof}
\begin{remark}\label{tor}
We henceforth regard $T^L$ as a subgroup of $G(O)$ in the unramified case and of $G'(O)$ in the ramified case.
\end{remark}
\subsection{Subgroups of $G(O)$ and $G'(O)$}
\begin{definition}\label{G-prime-filtration}
\begin{enumerate}
\item For $n\in\ZZ_{>0}$, let $G(O)(n)$ be the $n$-th congruence subgroup of $G(O)$.
\item Put $G(O)_n=G(O)/G(O)(n+1)$, and, for $r\leq n$, let $G(O)_n(r)$ be the image of $G(O)(r)$ in $G(O)_n$.
\item For $r\in\frac12\ZZ_{>0}$, let $G'(O)(r)$ be represented by matrices $\begin{pmatrix}a&b\\c&d\end{pmatrix}$ satisfying
$$
a-1,d-1\in\fm^{[r+1/2]},\qquad b\in\fm^{[r]},\qquad c\in\fm^{[r]+1},
$$
where $[r]$ is the integral part of $r$.
\item Put $G'(O)_r=G'(O)/G'(O)(r+1/2)$, and, for $s\leq r$, let $G'(O)_r(s)$ be the image of $G'(O)(s)$ in $G'(O)_r$.
\end{enumerate}
\end{definition}
\begin{remark}
The groups $G(O)_m$ and $G'(O)_m$ are algebraic $F$-groups.
\end{remark}
\begin{definition}
\begin{enumerate}
\item If $L/K$ is unramified, put $T^L(m)=T^L\cap G(O)(m)$.
\item If $L/K$ is ramified and $m>0$, put $T^L(m)=T^L\cap G'(O)(m-1/2)$.
\end{enumerate}
\end{definition}
\subsection{Local fields: notation}
From now on, $F$ is a local non-Archimedean field of odd residue characteristic.
\begin{definition}\label{ti-psi}
\begin{enumerate}
\item $O_F\subset F$ is the ring of integers.
\item $\varpi\in O_F$ is a uniformizer, $k=O_F/\varpi O_F$, and $q=|k|$.
\item $O_F^*(m)=1+\varpi^mO_F$.
\item $(O_F^*)_m=O_F^*/O_F^*(m+1)$.
\item For $m\geq n$, let $(O_F^*)_m(n)$ be the image of $O_F^*(n)$ in $(O_F^*)_m$.
\item Fix an additive character $\widetilde\psi:F\to\CC^*$ which is trivial on $\varpi O_F$ and nontrivial on $O_F$. Put $\widetilde\psi_n(x)=\widetilde\psi(\varpi^{-n}x)$.
\end{enumerate}
\end{definition}
\begin{remark}
The character $\widetilde\psi_n$ is trivial on $\varpi^{n+1}O_F$ and nontrivial on $\varpi^nO_F$. The set of such characters is a transitive $O_F^*$-space.
\end{remark}
\begin{definition}\label{psi}
\begin{enumerate}
\item Define $\psi:K\to\CC^*$ by
$$
\psi(f)=\widetilde\psi\left(\operatorname{Res}_{t=0}(f\,dt)\right).
$$
\item For $b\in K$, put $\psi_b(f)=\psi(bf)$.
\item For $n\in\ZZ$, put $\psi_n=\psi_{t^{-n-1}}$. Thus $\psi=\psi_{-1}$.
\end{enumerate}
\end{definition}
\begin{definition}\label{qext1}
Let $L/K$ be quadratic.
\begin{enumerate}
\item Let $\Xi^L$ be the set of smooth characters $\theta$ of $T^L$ such that $\theta\neq\overline\theta$.
\item Let $r$ be the least integer such that $\theta|_{T^L(r+1)}=1$. Define the depth $m(\theta)$ to be $r$ if $L/K$ is unramified and $r-\frac{1}{2}$ if $L/K$ is ramified.
\item Put $\Xi^L(m)=\{\theta\in\Xi^L:m(\theta)=m\}$.
\item Let $\bar\Xi^L(m)$ be the quotient of $\Xi^L(m)$ by $\theta\mapsto\theta^{-1}$.
\item Let $\ti\Xi^L(m)$ be the quotient of $\bar\Xi^L(m)$ by multiplication by quadratic characters of $T^L$.
\item Put $\D_L=\Hom(T^L,\ZZ/2\ZZ)$ and $\D_E=\Hom(T^E,\ZZ/2\ZZ)$.
\end{enumerate}
\end{definition}
\begin{lemma}\label{tori}
\begin{enumerate}
\item If $L/K$ is ramified, then $\D_L\simeq\ZZ/2\ZZ$.
\item If $L=E((t))$, then $\D_L\simeq\D_E$.
\item If $E/F$ is ramified, then $\D_E\simeq\ZZ/2\ZZ$.
\end{enumerate}
\end{lemma}
\begin{proof}
Every homomorphism to $\ZZ/2\ZZ$ factors through the quotient by squares. Since $1+tO$ is uniquely $2$-divisible by Hensel's lemma,
$$
K^*/(K^*)^2\simeq\langle t\rangle/\langle t^2\rangle\times F^*/(F^*)^2.
$$
If $L/K$ is ramified, write $L=K(x)$ with $x^2\in K$ of valuation $1$. The unit quotient contributes no additional square class, and
$$
L^*/K^*(L^*)^2\simeq\ZZ/2\ZZ,
$$
which proves (1). If $L=E((t))$, then
$$
L^*/K^*(L^*)^2\simeq E^*/F^*(E^*)^2,
$$
which proves (2). Applying the same argument to a ramified quadratic extension $E/F$ proves (3).
\end{proof}
\subsection{Schwartz spaces of $F$-varieties}
Let $Y$ be the set of $F$-points of an algebraic variety, with its natural topology. Let $\mcS(Y)$ be the space of locally constant complex-valued functions with compact support on $Y$. If $U\subset Y$ is open and $Z=Y\sm U$, then extension by zero and restriction give an exact sequence
$$
0\longrightarrow\mcS(U)\longrightarrow\mcS(Y)\longrightarrow\mcS(Z)\longrightarrow0.
$$
\section{Some results from classical representation theory}\label{reptheory}
All groups in this section are algebraic $F$-groups, all homomorphisms are algebraic, and all representations are complex and smooth.
\subsection{Coinvariants}
\begin{definition}\label{Vchi}
\begin{enumerate}
\item A linear algebraic group $H$ over $F$ is \emph{almost unipotent} if $H/R_u(H)$ is $F$-anisotropic, equivalently if $(H/R_u(H))(F)$ is compact.
\item A representation $\bs:H\to\Aut(V)$ is smooth if the stabilizer of every $v\in V$ is open.
\item Let $\mcR(H)$ denote the category of smooth representations of $H$.
\item Let $\widehat H$ denote the set of isomorphism classes of smooth irreducible representations of $H$. So if $H$ is commutative then  $\widehat H$ denotes the set of continuous characters of $H$.
\item If $(\bs,V)$ is a representation of $H$ and $\bl$ is a character of $H$, let $V(\bl)$ be the span of
$$
\{\bs(h)v-\bl(h)v:h\in H,\ v\in V\},
$$
and put $V_\bl=V/V(\bl)$.
\item If $H$ is commutative, put
$$
S(\bs)=\{\bl\in\widehat H:V_\bl\neq0\},
$$
and call it the support of $\bs$.
\end{enumerate}
\end{definition}
\begin{claim}
If $H$ is almost unipotent, then the functor $V\mapsto V_\bl$ from $\mcR(H)$ to $\Vect$ is exact for every character $\bl$ of $H$.
\end{claim}
\subsection{Compact induction}
\begin{definition}\label{indu}
Let $H$ be an algebraic $F$-group, let $L\subset H$ be closed, and let $(\bs,V)$ be a smooth representation of $L$. The compact induction $\ind_L^H(\bs)$ is the space of locally constant functions $f:H\to V$ satisfying
\begin{enumerate}
\item $f(hl^{-1})=\bs(l)f(h)$ for $h\in H$ and $l\in L$;
\item $\supp(f)\subset CL$ for some compact subset $C\subset H$.
\end{enumerate}
The group $H$ acts by left translation.
\end{definition}
Equivalently, let $\pi:H\to H/L$ be the quotient map and let $\mcF_\bs$ be the locally constant sheaf associated with the homogeneous bundle $H\times^L V$. Then $\ind_L^H(\bs)$ is the space of compactly supported sections of $\mcF_\bs$.
\begin{proposition}\label{indHL}
Let $H$, $L$, and $\bs$ be as above. Let $H'\subset H$ be a closed subgroup such that $H'L=H$, and put $L'=H'\cap L$. Then
$$
\operatorname{Res}_{H'}^H\ind_L^H(\bs)\simeq\ind_{L'}^{H'}(\bs|_{L'}).
$$
\end{proposition}
\begin{proof}
The natural map $H'/L'\to H/L$ is a continuous bijection. Both spaces are locally compact and Hausdorff. It remains to show that the map is open. The group $H'$ is a $\sigma$-compact, locally compact, totally disconnected group, and it has a basis of identity neighbourhoods consisting of compact open subgroups. The assertion therefore follows from Lemma \ref{open} below.


\begin{lemma}\label{open}
Let $A$ be a topological group acting transitively on a topological space $X, x\in X $, and  $\St_A(x)$ be the stabilizer of $x$ in $A$. Assume that
\begin{enumerate}
\item $X$ is locally compact and Hausdorff;
\item $A$ is $\sigma$-compact and the identity has a basis of neighbourhoods consisting of compact open subgroups.
\end{enumerate}
Then the orbit map $f:A\to X$, $a\mapsto a(x)$, is open. Equivalently, $A/\St_A(x)\to X$ is a homeomorphism.
\end{lemma}
\begin{proof}
It suffices to prove that $f(K)$ is open for every compact open subgroup $K\subset A$. Since $A$ is $\sigma$-compact and $K$ is open, the quotient $A/K$ is a $\sigma$-compact discrete space. Every compact subset of a discrete space is finite, so $A/K$ is countable. Choose representatives $a_1,a_2,\ldots$ for $A/K$. Then
$$
X=\bigcup_i f(a_iK),
$$
and every $f(a_iK)$ is compact. Since a locally compact Hausdorff space is a Baire space, one of these compact subsets has nonempty interior. As $f(a_iK)=a_i f(K)$, the set $f(K)$ has nonempty interior. The group $K$ acts transitively on $f(K)$, so every point of $f(K)$ is an interior point. Thus $f(K)$ is open.
\end{proof}
\end{proof}
\subsection{Induction from a character}
\begin{definition}
For a character $\bl:L\to\CC^*$, define
$$
\kk_\bl:\mcS(H)\longrightarrow\mcS(H)_{L,\bl},\qquad
\kk_\bl(f)(h)=\int_L\bl(l)f(hl)\,dl,
$$
where $dl$ is a Haar measure on $L$.
\end{definition}
\begin{claim}
The map $\kk_\bl$ identifies $\ind_L^H(\bl)$ with the representation of $H$ on the right $(L,\bl)$-coinvariants $\mcS(H)_{L,\bl}$, with $H$ acting by left translation.
\end{claim}
Let $Y$ be an $F$-variety with an action of $L$, and let $\bl$ be a character of $L$. For $y\in Y$, write $\St_L(y)$ for its stabilizer and put
$$
Y(\bl)=\{y\in Y:\bl|_{\St_L(y)}=1\}.
$$
\begin{lemma}\label{relevant-orbit}
Let $V(\bl)=\mcS(Y)_{L,\bl}$.
\begin{enumerate}
\item If $Y(\bl)=\emp$, then $V(\bl)=0$.
\item If $L$ is unipotent, $Y(\bl)\neq\emp$, and $L$ acts transitively on $Y(\bl)$, then $\dim V(\bl)=1$.
\end{enumerate}
\end{lemma}
\begin{proof}
The first assertion follows from \cite[Section~6, Theorems~6.9 and A]{BZ}. For (2), put $X=Y(\bl)$. It is the unique relevant $L$-orbit and is locally closed. Since $(L,\bl)$-coinvariants are exact, the exact sequences associated with
$$
X\subset\overline X\subset Y
$$
show that
$$
\mcS(Y)_{L,\bl}\simeq\mcS(\overline X)_{L,\bl}\simeq\mcS(X)_{L,\bl}.
$$
The last space is one-dimensional because $L$ acts transitively on $X$ and $\bl$ is trivial on the stabilizer of a point of $X$.
\end{proof}
\subsection{Mackey theory}
\begin{definition}\label{Mac0}
Let $H$ be an algebraic $F$-group and let $U\subset H$ be a closed normal commutative unipotent subgroup.
\begin{enumerate}
\item Let $Z(H)$ be the center of $H$.
\item For $\pi\in\widehat H$, put $S_\pi(U)=S(\pi|_U)$.
\item For $\bl\in\widehat U$, let $H_\bl=\St_H(\bl)$, and let $\widetilde H_\bl\subset\widehat{H_\bl}$ consist of the irreducible representations $(\rho,V)$ such that
$$
\rho(u)=\bl(u)\Id_V\qquad(u\in U).
$$
\item If $(\pi,V)\in\widehat H$ and $\bl\in S_\pi(U)$, let $\rho(\bl,\pi)$ be the natural representation of $H_\bl$ on $V_\bl$.
\item If $\rho\in\widetilde H_\bl$, put $\pi(\bl,\rho)=\ind_{H_\bl}^H(\rho)$.
\end{enumerate}
\end{definition}
\begin{lemma}\label{MM}
Let $\pi$ be an irreducible representation of $H$.
\begin{enumerate}
\item The set $S_\pi(U)$ is $\Ad(H)$-invariant, and $H$ acts transitively on it.
\item For every $\bl\in\widehat U$ and every $\rho\in\widetilde H_\bl$, the representation $\pi(\bl,\rho)$ is irreducible.
\item For every $\bl\in S_\pi(U)$, the representation $\rho(\bl,\pi)$ belongs to $\widetilde H_\bl$ and is irreducible.
\item For $\rho\in\widetilde H_\bl$, one has $\rho(\bl,\pi(\bl,\rho))\simeq\rho$.
\end{enumerate}
\end{lemma}
\begin{proof}
We reduce the statement to the usual Mackey theorem for a semidirect product. Let
$$
\widetilde H=H\ltimes U
$$
with multiplication
$$
(h,u)(h',u')=(hh',h'^{-1}uh'\,u').
$$
Then
$$
q:\widetilde H\longrightarrow H,\qquad q(h,u)=hu,
$$
is a surjective homomorphism, with
$$
Q=\ker(q)=\{(u^{-1},u):u\in U\}.
$$
Inflation along $q$ is a bijection between representations of $H$ and representations of $\widetilde H$ which are trivial on $Q$.
Apply \cite[Theorem~2]{H} to the normal abelian subgroup $\{1\}\ltimes U$ of $\widetilde H$. Its character orbit is the ordinary $H$-orbit in $\widehat U$, and the stabilizer of $\bl$ is $H_\bl\ltimes U$. If $\rho\in\widetilde H_\bl$, define
$$
\widetilde\rho(h,u)=\rho(h)\bl(u),\qquad h\in H_\bl,\ u\in U.
$$
Because $H_\bl$ stabilizes $\bl$, this is a representation of $H_\bl\ltimes U$. Moreover,
$$
\widetilde\rho(u^{-1},u)=\rho(u^{-1})\bl(u)=1,
$$
so it is trivial on $Q\cap(H_\bl\ltimes U)$. The induced representation of $\widetilde H$ is therefore trivial on $Q$ and descends to $H$; under the quotient map $q$, the descended representation is precisely $\ind_{H_\bl}^H(\rho)$. Conversely, the little-group representation attached by Mackey theory to an irreducible representation of $\widetilde H$ trivial on $Q$ is itself trivial on the relevant intersection with $Q$, and hence descends to a representation in $\widetilde H_\bl$. The four assertions now follow from the corresponding assertions of Mackey theory.
\end{proof}
\subsubsection{Irreducibility of restrictions}
Let $H$ be an algebraic $F$-group, let $U\subset H$ be a closed normal commutative unipotent subgroup, let $\lambda\in\widehat U$, and put $H_\lambda=\St_H(\lambda)$. For a closed subgroup $H'\subset H$, put
$$
U'=U\cap H',\qquad \lambda'=\lambda|_{U'},\qquad K=H'\cap H_\lambda,
$$
and let $\Omega=H\lambda\subset\widehat U$. For $\rho\in\widetilde H_\lambda$, put
$$
\pi(\lambda,\rho)=\ind_{H_\lambda}^H(\rho).
$$
\begin{lemma}\label{3.12}
Assume that the restriction map
$$
r:\Omega\longrightarrow\widehat{U'}
$$
is injective. Then $\pi(\lambda,\rho)|_{H'}$ is irreducible if and only if
$$
H'H_\lambda=H
$$
and $\rho|_K$ is irreducible.
\end{lemma}
\begin{proof}
Put $\Pi=\pi(\lambda,\rho)$. In the homogeneous-bundle realization of $\Pi$, the group $U$ acts on the fiber over $hH_\lambda$ by the character $h\lambda$. Hence
$$
S_\Pi(U)=\Omega,
$$
and after restriction to $U'$ one has
$$
S_{\Pi|_{H'}}(U')=r(\Omega).
$$
Suppose that $\Pi|_{H'}$ is irreducible. By Lemma \ref{MM}(1), $H'$ acts transitively on $r(\Omega)$. Since $r$ is $H'$-equivariant and injective, $H'$ acts transitively on $\Omega$. Thus $H'\lambda=H\lambda$, equivalently $H'H_\lambda=H$.
We claim that $\St_{H'}(\lambda')=K$. The inclusion $K\subset\St_{H'}(\lambda')$ is immediate. Conversely, if $h'\in H'$ stabilizes $\lambda'$, then
$$
r(h'\lambda)=h'r(\lambda)=r(\lambda),
$$
and injectivity of $r$ gives $h'\lambda=\lambda$. Hence $h'\in K$.
By Proposition \ref{indHL},
$$
\Pi|_{H'}\simeq\ind_K^{H'}(\rho|_K).
$$
If $\rho|_K$ were reducible, a nonzero proper $K$-subrepresentation would induce to a nonzero proper $H'$-subrepresentation. Here compact induction is exact and faithful: a nonzero vector gives a nonzero compactly supported local section of the associated homogeneous bundle. Thus $\rho|_K$ is irreducible.
Conversely, assume that $H'H_\lambda=H$ and that $\rho|_K$ is irreducible. Then $K=\St_{H'}(\lambda')$, and $U'$ acts on $\rho|_K$ by $\lambda'$. Proposition \ref{indHL} gives
$$
\Pi|_{H'}\simeq\ind_K^{H'}(\rho|_K),
$$
which is irreducible by Lemma \ref{MM}(2), applied to $(H',U')$.
\end{proof}
\subsection{Heisenberg groups}
\begin{definition}\label{dh}
Let $(W,\langle\ ,\ \rangle)$ be a finite-dimensional symplectic $F$-vector space, and put $\Sp(W)=\Sp(W,\langle\ ,\ \rangle)$.
\begin{enumerate}
\item Let $R_W=W\oplus F$ with multiplication
$$
(w',a')(w'',a'')=\left(w'+w'',a'+a''+\frac12\langle w',w''\rangle\right).
$$
We identify $Z(R_W)$ with $F$ by $a\mapsto(0,a)$. The group $\Sp(W)$ acts by
$$
(w,a)^g=(g(w),a).
$$
\item Let $\widehat R_W^{\widetilde\psi}$ be the set of isomorphism classes of irreducible smooth representations $(\pi,N)$ of $R_W$ such that
$$
\pi(0,a)=\widetilde\psi(a)\Id_N.
$$
\end{enumerate}
\end{definition}
\begin{claim}\label{hei}
\begin{enumerate}
\item There is, up to isomorphism, a unique irreducible smooth representation
$$
\tau_W:R_W\longrightarrow\Aut(V_W)
$$
with central character $\widetilde\psi$.
\item If $N$ is any smooth representation of $R_W$ on which $Z(R_W)$ acts by $\widetilde\psi$, then the evaluation map
$$
\operatorname{ev}_N:\Hom_{R_W}(V_W,N)\otimes V_W\longrightarrow N,
\qquad T\otimes v\longmapsto T(v),
$$
is an isomorphism.
\item There is a canonical central extension
$$
1\longrightarrow\{\pm1\}\longrightarrow\widetilde{\Sp(W)}\stackrel{\nu}{\longrightarrow}\Sp(W)\longrightarrow1
$$
and a metaplectic representation
$$
\kk_W:\widetilde{\Sp(W)}\longrightarrow\Aut(V_W)
$$
such that
$$
\tau_W(r^g)=\kk_W(\widetilde g)\tau_W(r)\kk_W(\widetilde g)^{-1}
$$
for $g\in\Sp(W)$, $r\in R_W$, and $\widetilde g\in\nu^{-1}(g)$.
\end{enumerate}
\end{claim}
\begin{proof}
Parts (1) and (3) are the Stone--von Neumann theorem and the construction of the Weil representation; see \cite{W}. We prove (2). Choose a polarization $W=X\oplus Y$. In the Schrödinger model, $V_W=\mcS(X)$, and
$$
\tau_W(x_0)f(x)=f(x+x_0),\qquad
\tau_W(y_0)f(x)=\widetilde\psi(\langle x,y_0\rangle)f(x).
$$
Let $\mcH_{\widetilde\psi}(R_W)$ be the twisted Hecke algebra of locally constant functions on $R_W$ which are compactly supported modulo $Z(R_W)$ and satisfy
$$
f(zr)=\widetilde\psi(z)^{-1}f(r).
$$
Its action on $\mcS(X)$ identifies it with the algebra of finite-rank operators on $\mcS(X)$; more precisely, it is a directed union of full matrix algebras obtained by imposing compact-support and compact-open-invariance conditions on $X$. Every nondegenerate module over a full matrix algebra is a direct sum of copies of its defining module, with multiplicity space the Hom-space from that module. Passing through the directed union gives
$$
N\simeq\Hom_{R_W}(V_W,N)\otimes V_W,
$$
and the resulting isomorphism is the evaluation map.
\end{proof}
\begin{definition}\label{M}
Let $M\subset\Sp(W)$ be a closed subgroup. A splitting of $\nu$ over $M$ is a continuous homomorphism
$$
s:M\longrightarrow\widetilde{\Sp(W)}
$$
such that $\nu\circ s=\Id_M$. Let $\mcS_M$ be the set of splittings.
\end{definition}
\begin{claim}\label{Z2}
If $\mcS_M\neq\emp$, then $\mcS_M$ is a torsor under $\Hom_{\mathrm{cont}}(M,\{\pm1\})$.
\end{claim}
\begin{proof}
If $s_1$ and $s_2$ are splittings, then $m\mapsto s_1(m)s_2(m)^{-1}$ is a continuous homomorphism to the central kernel. Conversely, multiplying a splitting by such a homomorphism gives another splitting.
\end{proof}
\begin{claim}\label{heis}
The metaplectic double cover splits over every compact subgroup $M\subset\Sp(W)$.
\end{claim}
\begin{proof}
The orbit of a point in the Bruhat--Tits building under $M$ is bounded, so the Bruhat--Tits fixed-point theorem gives a fixed point. Hence $M$ is contained in the stabilizer of a vertex and therefore in a maximal parahoric subgroup. By \cite[Proposition~4.1]{HW}, the metaplectic cover splits over every maximal parahoric subgroup of $\Sp(W)$ in odd residue characteristic. Restriction gives a splitting over $M$.
\end{proof}
\subsection{The two-dimensional case}
Assume that $\dim_F(W)=2$. A nonzero alternating form on $W$ is unique up to a scalar, and $\Sp(W)\simeq\SL_2(F)$.
\begin{lemma}\label{unique-splitting}
The metaplectic cover has a unique splitting over $\SL_2(O_F)$.
\end{lemma}
\begin{proof}
Existence follows from Claim \ref{heis}. The ratio of two splittings is a continuous homomorphism
$$
\SL_2(O_F)\longrightarrow\{\pm1\}.
$$
Let $K_1$ be the kernel of reduction to $\SL_2(k)$. Since $K_1$ is pro-$p$ and $p\neq2$, the homomorphism is trivial on $K_1$ and factors through $\SL_2(k)$. The latter group is generated by its two root subgroups, each a finite $p$-group, so every homomorphism $\SL_2(k)\to\{\pm1\}$ is trivial.
\end{proof}
Let $E/F$ be quadratic. Multiplication on the two-dimensional $F$-vector space $E$ embeds
$$
E^1=\{u\in E^*:N_{E/F}(u)=1\}
$$
into $\SL_F(E)\simeq\SL_2(F)$. Let $\mcC(E)$ be the set of maximal compact subgroups of $\SL_F(E)$ containing $E^1$.
\begin{lemma}\label{lif}
\begin{enumerate}
\item If $E/F$ is unramified, then $|\mcC(E)|=1$.
\item If $E/F$ is ramified, then $\mcC(E)=\{C_+,C_-\}$, and $C_+\cap C_-$ is an Iwahori subgroup.
\end{enumerate}
\end{lemma}
\begin{proof}
Vertices of the Bruhat--Tits tree of $\SL_2(F)$ are homothety classes of $O_F$-lattices in $E$. If the class of a lattice $\Lambda$ is fixed by $E^1$, then for every $u\in E^1$ one has $u\Lambda=a_u\Lambda$ for some $a_u\in F^*$. Comparison of covolumes and $N_{E/F}(u)=1$ give $v_F(a_u)=0$, so $a_u\Lambda=\Lambda$ and $u\Lambda=\Lambda$.
Choose $u\in E^1$ with $O_F[u]=O_E$. In the ramified case, write $E=F(\vartheta)$ with $\vartheta^2=\varpi\epsilon$ and take
$$
u=\frac{1+\vartheta}{1-\vartheta}.
$$
Then $u\bar u=1$ and $\vartheta=(u-1)/(u+1)$, with $u+1$ a unit. In the unramified case, choose a norm-one unit whose reduction generates the residue extension. Thus every $E^1$-stable lattice is an $O_E$-module and hence a fractional ideal $\mathfrak p_E^r$. Two such ideals are $F^*$-homothetic precisely when
$$
r\equiv s\pmod{e(E/F)}.
$$
There is therefore one fixed vertex in the unramified case and two adjacent fixed vertices in the ramified case. Their stabilizers are the asserted maximal compact subgroups, and the stabilizer of the connecting edge is their Iwahori intersection. See \cite{S}.
\end{proof}
\begin{corollary}\label{lift}
If $E/F$ is unramified, then the torsor $\mcS_{E^1}$ has a canonical trivialization.
\end{corollary}
\begin{proof}
By Lemma \ref{lif}, $E^1$ lies in a unique maximal compact subgroup $C$. The metaplectic cover has a unique splitting over $C$ by Lemma \ref{unique-splitting}, after conjugating $C$ to $\SL_2(O_F)$. Its restriction to $E^1$ is canonical.
\end{proof}
Assume now that $E/F$ is ramified. Let $C_+$ and $C_-$ be as in Lemma \ref{lif}, and put $I=C_+\cap C_-$. Choose $b\in E^*$ of odd $E$-valuation and let
$$
\beta_b(g)=bgb^{-1},\qquad g\in\SL_F(E),
$$
where $b$ acts on $E$ by multiplication.
\begin{claim}\label{beta-swap}
One has $\beta_b(C_+)=C_-$ and $\beta_b(C_-)=C_+$.
\end{claim}
\begin{proof}
Multiplication by $b$ sends $\mathfrak p_E^r$ to $\mathfrak p_E^{r+v_E(b)}$. Since $v_E(b)$ is odd, it interchanges the two parity classes of fixed vertices and hence their stabilizers.
\end{proof}
\begin{claim}\label{beta-lift}
There is a unique automorphism
$$
\widetilde\beta_b:\widetilde{\SL_2(F)}\longrightarrow\widetilde{\SL_2(F)}
$$
which fixes $\{\pm1\}$ and satisfies $\nu\circ\widetilde\beta_b=\beta_b\circ\nu$.
\end{claim}


\begin{proof}
Pullback by $\beta_b$ preserves the metaplectic cohomology class; this follows from Moore's classification or directly from the cocycle of Ranga Rao, so a lift exists. The ratio of two lifts is a continuous homomorphism $\SL_2(F)\to\{\pm1\}$. Each root subgroup is isomorphic to $(F,+)$ and is uniquely $2$-divisible, so this homomorphism is trivial on both root subgroups and therefore on $\SL_2(F)$. This proves uniqueness. Multiplicativity follows from uniqueness because both sides lift conjugation by $b_1b_2$; see \cite{Moore,RangaRao,Roberts}.
\end{proof}
Let $s_\pm:C_\pm\to\widetilde{\SL_2(F)}$ be the unique splittings. Define
$$
\chi_I(i)=s_+(i)s_-(i)^{-1},\qquad i\in I,
$$
and put $\delta=\chi_I|_{E^1}$.
\begin{lemma}\label{ex}
The character $\chi_I$ is the unique nontrivial quadratic character of the reductive quotient $I/I^+\simeq k^*$. With the standard normalization, if $i\in I$ has diagonal reduction $\operatorname{diag}(\bar a,\bar a^{-1})$ and $\eta:k^*\to\{\pm1\}$ is the quadratic character, then
$$
\chi_I(i)=\eta(\bar a).
$$
Consequently,
$$
\delta(-1)=\eta(-1)=(-1)^{(q-1)/2},
$$
and
$$
\delta=\begin{cases}
1,&q\equiv1\pmod4,\\
\text{the unique nontrivial quadratic character of }E^1,&q\equiv3\pmod4.
\end{cases}
$$
\end{lemma}
\begin{proof}

We first observe that $\chi_I$ is a continuous homomorphism: its values
$s_+(i)s_-(i)^{-1}$ lie in the central kernel $\{\pm1\}$ of $\nu$, so
for $i_1,i_2\in I$ we have
\[
\chi_I(i_1 i_2)=s_+(i_1)s_+(i_2)\,s_-(i_2)^{-1}s_-(i_1)^{-1}
=\chi_I(i_1)\chi_I(i_2),
\]
the central factor $\chi_I(i_2)$ having been moved across
$s_-(i_1)^{-1}$; continuity follows from that of $s_\pm$.

The action on the Bruhat--Tits tree gives the amalgam decomposition
\[
\mathrm{SL}_2(F)=C_+\ast_{I}C_-.
\]
If $\chi_I$ were trivial, the splittings $s_+$ and $s_-$ would agree on
$I$ and hence would glue to a splitting of the metaplectic cover over
$\mathrm{SL}_2(F)$, contradicting the nonsplitting of the standard
metaplectic cover. Hence $\chi_I\neq 1$. Since $\chi_I$ takes values in
$\{\pm1\}$ and $I^+$ is pro-$p$ with $p\neq 2$, it is trivial on $I^+$
and factors through a quadratic character of $I/I^+\simeq k^*$. The
group $k^*$ is cyclic of even order $q-1$, so it has exactly one
nontrivial quadratic character, namely $\eta$. Thus $\chi_I$ is the
composition of the reduction $i\mapsto\bar a$ with $\eta$, which is the
displayed formula; it agrees with the value given by the explicit
metaplectic cocycle in its standard normalization \cite{Moore,
RangaRao}.

We now compute $\delta=\chi_I|_{E^1}$. Write $E=F(\vartheta)$ with
$\vartheta^2=\varpi\epsilon$, $\epsilon\in\mathcal{O}_F^*$. In the
basis $(\vartheta,1)$ we have
$\mathcal{O}_E=\mathcal{O}_F\vartheta\oplus\mathcal{O}_F$ and
$\mathfrak{p}_E=\vartheta\mathcal{O}_E=
\mathcal{O}_F\vartheta\oplus\varpi\epsilon\,\mathcal{O}_F$, so this
basis identifies the lattice chain
$\mathcal{O}_E\supset\mathfrak{p}_E$ with the standard chain, and the
groups $C_\pm$ and $I$ become the standard maximal compact subgroups
and the standard Iwahori subgroup. Multiplication by $a+b\vartheta$ is
represented by the matrix
\[
\begin{pmatrix} a & b\\ b\varpi\epsilon & a\end{pmatrix}.
\]
If $a+b\vartheta\in E^1$, then $a^2-b^2\varpi\epsilon=1$; reduction
modulo $\varpi$ gives $\bar a^2=1$, so $\bar a=\pm 1$. In particular
$a\in\mathcal{O}_F^*$, the matrix lies in $I$, and its diagonal
reduction $\operatorname{diag}(\bar a,\bar a)$ equals
$\operatorname{diag}(\bar a,\bar a^{-1})$. Hence
$\delta(a+b\vartheta)=\eta(\bar a)$. The element $-1\in E^1$ reduces to
$\operatorname{diag}(-1,-1)$, so the image of $E^1$ in
$I/I^+\simeq k^*$ is exactly $\{\pm1\}$, and $\delta$ is nontrivial if
and only if $\eta(-1)=-1$.

Finally, the kernel of the reduction map $E^1\to\{\pm1\}$ is pro-$p$,
so $E^1$ has at most one nontrivial quadratic character. Since
$\eta(-1)=(-1)^{(q-1)/2}$, evaluation at $-1$ gives the asserted
dichotomy.
\end{proof}
\begin{corollary}\label{no}If the extension $E/F$ is ramified there is no $\widetilde\beta_b$-equivariant splitting of the metaplectic cover over $G'(O_F)$ (see Subsection \ref{classical}).
\end{corollary}


\subsection{Quasi-Heisenberg groups}
\begin{definition}\label{qH}
Let $W$ be symplectic. A quasi-Heisenberg group is a quintuple
$$
(H,R_W,S,\widetilde a,\bl)
$$
with the following properties.
\begin{enumerate}
\item $H$ is an algebraic $F$-group containing $R_W$ as a closed normal subgroup, and $Z(R_W)\subset Z(H)$.
\item $S$ is a closed commutative subgroup of $H$ such that multiplication gives a bijection
$$
S\times Z(H)R_W\longrightarrow H.
$$
Equivalently, $S\to H/Z(H)R_W$ is an isomorphism.
\item The homomorphism $\widetilde a:H\to\widetilde{\Sp(W)}$ lifts the action $a:H\to\Sp(W)$ induced by conjugation on $W=R_W/Z(R_W)$.
\item The character $\bl\in\widehat{Z(H)}$ restricts to $\widetilde\psi$ on $Z(R_W)$.
\end{enumerate}
Let $\operatorname{pr}_S:H\to S$ be the projection induced by the displayed product decomposition.
\end{definition}
Let $\widetilde{\widehat H}_\bl$ be the set of irreducible smooth representations $(\rho,N)$ of $H$ with central character $\bl$.
\begin{claim}\label{qH-phi}
For a quasi-Heisenberg group there is a unique representation
$$
\phi:H\longrightarrow\Aut(V_W)
$$
such that
$$
\phi|_{R_W}=\tau_W,\qquad \phi|_{Z(H)}=\bl\Id_{V_W},\qquad \phi|_S=\kk_W\circ\widetilde a.
$$
\end{claim}
\begin{proof}
The prescriptions agree on pairwise intersections. The covariance identity in Claim \ref{hei}(3) gives compatibility with the conjugation action of $S$ on $R_W$. Since $S$, $Z(H)$, and $R_W$ generate $H$, the representation exists and is unique.
\end{proof}
\begin{lemma}\label{thet}
Let $H$ be quasi-Heisenberg and let $\rho\in\widetilde{\widehat H}_\bl$. Then
\begin{enumerate}
\item $\rho|_{R_W}$ is irreducible;
\item there is a unique character $\theta_\rho$ of $S$ such that
$$
\rho=(\theta_\rho\circ\operatorname{pr}_S)\otimes\phi.
$$
\end{enumerate}
\end{lemma}
\begin{proof}
By Claim \ref{hei}(2), evaluation gives
$$
N\simeq\overline N\otimes V_W,
\qquad \overline N=\Hom_{R_W}(V_W,N).
$$
The representations $\rho$ and $\phi$ define an action of $H$ on $\overline N$ by
$$
(h\cdot T)=\rho(h)T\phi(h)^{-1}.
$$
This action is trivial on $R_W$ and on $Z(H)$, so it factors through $H/Z(H)R_W\simeq S$. Under the displayed tensor decomposition, $H$-stable subspaces correspond to $S$-stable subspaces of $\overline N$. Since $\rho$ is irreducible and $S$ is commutative, $\overline N$ is one-dimensional. Its character is $\theta_\rho$, and the one-dimensionality also proves (1).
\end{proof}
Define
$$
\alpha_{\widetilde a}:\widehat S\longrightarrow\widetilde{\widehat H}_\bl,
\qquad
\alpha_{\widetilde a}(\theta)=(\theta\circ\operatorname{pr}_S)\otimes\phi.
$$
\begin{corollary}\label{alpha}
\begin{enumerate}
\item The map $\alpha_{\widetilde a}$ is a bijection.
\item If $\delta\in\Hom(S,\{\pm1\})$ and $\widetilde a$ is replaced by $\delta\widetilde a$, then
$$
\alpha_{\delta\widetilde a}(\theta)=\alpha_{\widetilde a}(\delta\theta).
$$
\end{enumerate}
\end{corollary}
\begin{proof}
Part (1) is Lemma \ref{thet}. For (2), changing the lift by $\delta$ changes $\phi$ by the one-dimensional factor $\delta\circ\operatorname{pr}_S$.
\end{proof}
\section{Special representations of $P(O)$}\label{PO}
\subsection{Classification of smooth irreducible representations of $P(O)$}
The group $P(O)=O^*\ltimes O$ is a pro-algebraic $F$-group, and every smooth irreducible representation factors through a finite-dimensional quotient. Mackey theory therefore applies. Every nontrivial continuous character of $O$ is conjugate under $O^*$ to $\psi_m$ for a unique $m\geq0$, where
$$
\psi_m(x)=\psi(t^{-m-1}x).
$$
Its stabilizer is $O^*(m+1)\ltimes O$.
\begin{definition}\label{ph}
Let $m\geq0$.
\begin{enumerate}
\item For a character $\chi$ of $O^*(m+1)$, let $\widehat\chi$ be the character of $\St_{\psi_m}=O^*(m+1)\ltimes O$ given by
$$
\widehat\chi(a,x)=\chi(a)\psi_m(x).
$$
\item Put $\bs(m,\chi)=\ind_{\St_{\psi_m}}^{P(O)}(\widehat\chi)$.
\item Let $\mcS(m,\chi)$ be the representation space of $\bs(m,\chi)$.
\item Put $\bs(m)=\bs(m,1)$ and $\mcS(m)=\mcS(m,1)$.
\end{enumerate}
\end{definition}
\begin{lemma}\label{p-o}
\begin{enumerate}
\item The representations $\bs(m,\chi)$ are irreducible and pairwise nonisomorphic, and every irreducible representation of $P(O)$ of dimension greater than $1$ is of this form.
\item The subgroup $O^*(m+1)$ acts on $\bs(m,\chi)$ by the character $\chi$.
\end{enumerate}
\end{lemma}
\begin{proof}
This is the little-group classification of Lemma \ref{MM}, applied to the normal subgroup $O\subset P(O)$.
\end{proof}
\subsection{A second realization of $\bs(m)$}
Let $\widetilde R_m=O^*(1)\ltimes t^mO$. It has a character $\bl_m$ which is trivial on $O^*(1)$ and equals $\psi_m$ on $t^mO$.
\begin{lemma}\label{sec}
There is an isomorphism
$$
\bs(m)\simeq\ind_{\widetilde R_m}^{P(O)}(\bl_m).
$$
\end{lemma}
\begin{proof}
A direct Mackey calculation for the normal subgroup $O\subset P(O)$ shows that the $O$-support of $\ind_{\widetilde R_m}^{P(O)}(\bl_m)$ is the single orbit $O^*\psi_m$. Indeed, the characters of $O$ extending $\psi_m|_{t^mO}$ form the $O^*(1)$-orbit of $\psi_m$, and induction from $O^*(1)$ to $O^*$ produces the full $O^*$-orbit. The $\psi_m$-coinvariant space is one-dimensional, and $O^*(m+1)$ acts trivially on it. Lemma \ref{MM} therefore identifies the representation with $\bs(m)$.
\end{proof}
\subsection{Depth and special representations}
We use the notation of Definition \ref{G-prime-filtration}.
\begin{definition}\label{depth}
\begin{enumerate}
\item The depth $m(\pi)$ of a representation of $G(O)$ is the least $n\geq0$ such that $G(O)(n+1)$ acts trivially.
\item The depth $m(\pi)$ of a representation of $G'(O)$ is the least $n\in\frac12\ZZ_{\geq0}$ such that $G'(O)(n+1/2)$ acts trivially.
\item A representation of $G(O)$ or $G'(O)$ is special if it has dimension greater than $1$ and its restriction to $P(O)$ is irreducible.
\item An irreducible representation of $P(O)$ of dimension greater than $1$ is special if it extends to $G(O)$ or $G'(O)$.
\end{enumerate}
\end{definition}
\begin{proposition}\label{eq}
If $\bs(m,\chi)$ is special, then $\chi=1$.
\end{proposition}
\begin{remark}
We shall prove the converse below.
\end{remark}
\begin{proof}
Suppose first that $\bs(m,\chi)$ extends to a special representation $(\pi,V)$ of $G(O)$. Let $i:P(O)\ho G(O)$ be the standard embedding and choose a $P(O)$-equivariant isomorphism
$$
\kk_\pi:\mcS(m,\chi)\longrightarrow V.
$$
Let $s=\begin{pmatrix}0&1\\1&0\end{pmatrix}$, let $i'=\Ad(s)\circ i$, and put $\kk'_\pi=\pi(s)\circ\kk_\pi$. Then
$$
A(\pi)=\kk_\pi^{-1}\circ\kk'_\pi
$$
intertwines the two $P(O)$-actions. Since $i'(\widehat a)=i(\widehat{a^{-1}})$, one has
$$
A(\pi)\bs(m,\chi)(\widehat a)A(\pi)^{-1}=\bs(m,\chi)(\widehat{a^{-1}}).
$$
For $a\in O^*(m+1)$, Lemma \ref{p-o}(2) gives $\chi(a)=\chi(a^{-1})$, so $\chi^2=1$. The group $O^*(m+1)=1+t^{m+1}O$ is uniquely $2$-divisible because $2$ is invertible in $F$, and hence it has no nontrivial quadratic character. Thus $\chi=1$.
For $G'(O)$, use $s=\begin{pmatrix}0&1\\t&0\end{pmatrix}\in G'(O)$. Conjugation again sends $\widehat a$ to $\widehat{a^{-1}}$ in $\PGL(2)$, and the same argument proves $\chi=1$.
\end{proof}
\begin{lemma}\label{restr}
\begin{enumerate}
\item If $\pi$ is a special representation of $G(O)$ of depth $m$, then $\pi|_{P(O)}\simeq\bs(m)$.
\item If $\pi$ is a special representation of $G'(O)$ of depth $n-1/2$, where $n\in\ZZ_{>0}$, then $\pi|_{P(O)}\simeq\bs(n-1)$.
\end{enumerate}
\end{lemma}
\begin{proof}
By Proposition \ref{eq}, the restriction has the form $\bs(r)$. The representation $\bs(r)$ is trivial on the additive subgroup $t^{r+1}O$ and nontrivial on $t^rO$. Comparing this conductor with the defining congruence filtration gives $r=m$ in the first case and $r=n-1$ in the second.
\end{proof}
\subsection{The restriction of representations of $G(O)$ to $O\subset P(O)$}
We identify the additive group $O$ with $N(O)\subset P(O)$. Let $G^+(F)$ be the image of $\SL_2(F)$ in $\PGL_2(F)$, and let $G^+(O)$ be its inverse image in $G(O)$ under reduction modulo $t$. We use the analogous notation $G^+(O)_n$ for the image in $G(O)_n$.
\begin{proposition}\label{infty}
For every smooth representation $(\pi,V)$ of $G^+(O)$, the natural inclusion
$$
V^{G^+(O)}\hookrightarrow V^O
$$
is an isomorphism.
\end{proposition}
\begin{proof}
It is enough to prove the assertion for a representation which factors through $G^+(O)_n$. We argue by induction on $n$. For $n=0$, let $v$ be fixed by $N(F)$. By smoothness, $v$ is fixed by an open subgroup $J\subset G^+(F)$. Hence $v$ is fixed by $j^{-1}N(F)j$ for every $j\in J$. These conjugates generate the image of $\SL_2(F)$ in $\PGL_2(F)$, so $v$ is $G^+(F)$-fixed.
For the induction step, let $v$ be fixed by $O=N(O)$. The same argument applied to sufficiently small conjugates of $N(O)$ shows that $v$ is fixed by the last congruence subgroup $G(O)_n(n)\simeq\fg(F)$. Passing to $G^+(O)_{n-1}$ and applying the induction hypothesis shows that $v$ is fixed by $G^+(O)_n$.
\end{proof}
\section{Special representations of $G(O)$}\label{GO}
The purpose of this section is to classify and construct special irreducible representations of $G(O)$.
\subsection{Notation}
The group $G(O)_n$ acts by conjugation on the groups $G(O)_{2n-1}(n)$ and $G(O)_{2n}(n)$.
\begin{claim}
\begin{enumerate}
\item The group $G(O)_{2n-1}(n)$ is commutative.
\item For $n>0$, one has $Z(G(O)_{2n}(n))=G(O)_{2n}(n+1)$.
\item For $n>0$, the commutator subgroup of $G(O)_{2n}(n)$ is $G(O)_{2n}(2n)$.
\end{enumerate}
\end{claim}
\subsection{Lie algebras}
The trace pairing identifies $\fg^*$ with $\fg$. Let
$$
\iota:\fm\fg(O)\longrightarrow G(O)(1),\qquad X\longmapsto1+X,
$$
where the right-hand side is interpreted in $\PGL_2(O)$.
\begin{claim}\label{iota1}
\begin{enumerate}
\item The map $\iota$ is a bijection of sets. It induces bijections
$$
\iota_m(n):\fg(F)\otimes_F\fm^n/\fm^{m+1}\longrightarrow G(O)_m(n).
$$
\item If $m<2n$, then $\iota_m(n)$ is an isomorphism from the additive group on the left to the group on the right.
\end{enumerate}
\end{claim}
The map $\iota_m(m)$ identifies $G(O)_m(m)$ with $\fg(F)$ and, using $\widetilde\psi$, identifies its character group with $\fg(F)$. A character is called elliptic if the corresponding element of $\fg(F)$ is nonzero and elliptic.
Let
$$
H=G(O)_m,\qquad H'=P(O)_m,\qquad U=G(O)_m(m),
$$
and let $\bl\in\widehat U$ be elliptic.
\begin{claim}\label{P=G}
The restriction map from the $H$-orbit of $\bl$ to $\widehat{U\cap H'}$ is injective, and
$$
H'H_\bl=H.
$$
Thus the orbit and double-coset hypotheses in Lemma \ref{3.12} hold.
\end{claim}
\begin{proof}
Under the trace pairing, write
$$
Y=\begin{pmatrix}a&b\\c&-a\end{pmatrix}\in\fg(F).
$$
Restriction to the upper triangular Lie algebra records $(a,c)$, because
$$
\operatorname{tr}\left(Y\begin{pmatrix}x&y\\0&-x\end{pmatrix}\right)=2ax+cy.
$$
If $Y$ is elliptic, then $c\neq0$, and $b$ is determined by $(a,c)$ and the invariant $-\det(Y)=a^2+bc$. Hence restriction is injective on the orbit. The Borel group acts transitively on the corresponding set of pairs $(a,c)$, so it acts transitively on the elliptic orbit. This is equivalent to $H'H_\bl=H$.
\end{proof}
For a quadratic extension $E/F$, let $\ft^E$ be the Lie algebra of $T^E$, embedded in $\fg(F)$ through Claim \ref{imb}.
\begin{claim}\label{pro}
There is a unique $\Ad(T^E)$-equivariant projection $p:\fg(F)\to\ft^E$.
\end{claim}
The projection extends coefficientwise to $p(O):\fg(O)\to\ft^E(O)$. If $L=E((t))$, define
$$
p_m(n)=\iota_{T,m}(n)\circ p(O)\circ\iota_m(n)^{-1}:G(O)_m(n)\longrightarrow T^L_m(n),
$$
where $\iota_{T,m}(n)$ is the analogous coordinate map for the torus.
\begin{claim}\label{iota}
If $m<2n$, then $p_m(n)$ is a group homomorphism.
\end{claim}
Unless stated otherwise, we now consider irreducible representations of $G(O)$ of positive depth.
\subsection{Elliptic representations}
\begin{definition}\label{def}
Let $\pi$ be an irreducible representation of $G(O)$ of positive depth $m=m(\pi)$.
\begin{enumerate}
\item Let $\pi_m$ be the restriction of $\pi$ to $G(O)_m(m)$.
\item Regard $\pi_m$ as a representation of the additive group $\fg(F)$.
\item Let $\bo(\pi)\subset\fg^*(F)\simeq\fg(F)$ be its support.
\end{enumerate}
\end{definition}
\begin{lemma}\label{support}
The support $\bo(\pi)$ is a single coadjoint orbit.
\end{lemma}
\begin{proof}
This is Lemma \ref{MM}(1).
\end{proof}
\begin{definition}\label{def1}
\begin{enumerate}
\item The representation $\pi$ is elliptic if $\bo(\pi)$ is an elliptic orbit.
\item For elliptic $\pi$, let $E_\pi/F$ be the quadratic extension containing the eigenvalues of the elements of $\bo(\pi)$, and put $L_\pi=E_\pi((t))$.
\item For a quadratic extension $E/F$, let $\widehat G_m^E$ be the set of isomorphism classes of elliptic representations $\pi$ of depth $m$ with $E_\pi=E$.
\end{enumerate}
\end{definition}
\begin{lemma}\label{orbit-image}
Let $\mcO\subset\mathfrak{sl}_2(F)$ be a nonzero adjoint orbit, and let
$$
r:\mcO\longrightarrow\mathfrak p(F)^*
$$
be restriction to the upper triangular Lie algebra. Then $P(F)$ acts transitively on $r(\mcO)$ if and only if $\mcO$ is elliptic.
\end{lemma}
\begin{proof}
Write $Y=\begin{pmatrix}a&b\\c&-a\end{pmatrix}$ and put $q(Y)=a^2+bc=-\det(Y)$. Restriction to $\mathfrak p$ records $(a,c)$. If
$$
p=\begin{pmatrix}s&u\\0&1\end{pmatrix},
$$
then the induced action is
$$
(a,c)\longmapsto\left(a+\frac{u}{s}c,\frac{c}{s}\right).
$$
Thus the locus $c\neq0$ is one $P(F)$-orbit. If $q(Y)$ is nonsquare, then $c$ cannot vanish, so the orbit is elliptic and $P(F)$ is transitive on its restricted image. If $q(Y)=d^2\neq0$, the restricted image also contains the two points $(d,0)$ and $(-d,0)$, each a separate orbit. If $q(Y)=0$ and $Y\neq0$, the restricted image contains both the open orbit $c\neq0$ and the point $(0,0)$. Hence the restricted image is not transitive in the nonelliptic cases.
\end{proof}
\begin{lemma}\label{GO-special-elliptic}
Every special representation of $G(O)$ is elliptic.
\end{lemma}
\begin{proof}
The case $m=0$ is easy: in this case we are just dealing with an irreducible representation of $G(F)$ of dimension $>1$ and the fact that it is special means that its restriction to $P(F)$ is irreducible, which implies that it is cuspidal.

Let $m>0$ and let $\pi$ be an irreducible representation of $G(O)_m$. Put $U=G(O)_m(m)$ and $U'=U\cap P(O)_m$. The $U'$-support of $\pi|_{P(O)_m}$ is the image of $\bo(\pi)$ under restriction to $\mathfrak p(F)^*$. If $\pi|_{P(O)_m}$ were irreducible, Lemma \ref{MM}(1) would imply that $P(F)$ acts transitively on this image. Lemma \ref{orbit-image} then shows that $\bo(\pi)$ is elliptic.
\end{proof}
\subsection{Elliptic representations of odd depth}
Fix $n>0$ and put
$$
U=G(O)_{2n-1}(n),
$$
a commutative normal subgroup of $G(O)_{2n-1}$.
\begin{definition}\label{oddc}
Let $\bl\in\widehat U$.
\begin{enumerate}
\item Let $\bar\bl$ be the restriction of $\bl$ to $G(O)_{2n-1}(2n-1)\simeq\fg(F)$, regarded as an element of $\fg(F)$.
\item The character $\bl$ is elliptic if $\bar\bl$ is nonzero and elliptic.
\item For elliptic $\bl$, let $E_\bl/F$ contain the eigenvalues of $\bar\bl$, and put $L_\bl=E_\bl((t))$.
\item Let $\St_\bl$ be the stabilizer of $\bl$ in $G(O)_{2n-1}$.
\item Let $\bO_\bl$ be the $G(O)_{2n-1}$-orbit of $\bl$.
\item Let $\bO_\bl^{T^{L_\bl}}$ be the set of $T^{L_\bl}$-fixed points in $\bO_\bl$.
\end{enumerate}
\end{definition}
\begin{claim}\label{ellch}
Let $\bl$ be elliptic.
\begin{enumerate}
\item There is, up to conjugacy, a unique embedding $T^{L_\bl}\ho G(O)$ such that $T^{L_\bl}_{2n-1}\subset\St_\bl$.
\item One has $\St_\bl=T^{L_\bl}_{2n-1}U$.
\item One has $\bO_\bl^{T^{L_\bl}}=\{\bl,\bl^{-1}\}$.
\item One has $P(O)_{2n-1}\St_\bl=G(O)_{2n-1}$.
\end{enumerate}
\end{claim} We denote by $T^L_m$ the image of $T^L$ in $G(O)_m$ and by $T^L_m(r)$
the image of $T^L(r)$.


\begin{theorem}\label{charodd}
Every elliptic representation of $G(O)$ of odd depth is special.
\end{theorem}
\begin{proof}
Let $\pi$ be elliptic of depth $2n-1$. By Lemma \ref{MM}, there are an elliptic character $\bl_\pi\in\widehat U$, unique up to inversion among the $T^{L_\pi}$-fixed points of the support, and an irreducible representation $\rho_\pi$ of $\St_{\bl_\pi}$ such that
$$
\pi\simeq\ind_{\St_{\bl_\pi}}^{G(O)_{2n-1}}(\rho_\pi).
$$
Since $\St_{\bl_\pi}/\ker(\bl_\pi)$ is commutative, $\rho_\pi$ is a character. Claim \ref{ellch}(4), Claim \ref{P=G}, and Lemma \ref{3.12} show that $\pi|_{P(O)}$ is irreducible.
\end{proof}
Let $\theta_{\pi,\bl_\pi}$ be the restriction of $\rho_\pi$ to $T^{L_\pi}_{2n-1}$.
\begin{claim}\label{inv}
One has
$$
\theta_{\pi,\bl_\pi^{-1}}=\theta_{\pi,\bl_\pi}^{-1}.
$$
\end{claim}
\begin{proof}
The two fixed characters $\bl_\pi$ and $\bl_\pi^{-1}$ are exchanged by the nontrivial element of the normalizer of $T^{L_\pi}$. This element acts on $T^{L_\pi}$ by Galois conjugation, which is inversion under $T^{L_\pi}\simeq(L_\pi)^1$.
\end{proof}
Thus $\pi$ determines a point of $\bar\Xi^{L_\pi}(2n-1)$.
\subsection{Construction at odd depth}
Let $E/F$ be quadratic and put $L=E((t))$.
\begin{definition}
For $\theta\in\Xi^L(2n-1)$, define
$$
\bl_\theta=\theta\circ p_{2n-1}(n)\in\widehat U.
$$
\end{definition}
\begin{claim}\label{oo}
\begin{enumerate}
\item The character $\bl_\theta$ is elliptic.
\item There is a unique character $\widehat\theta$ of $\St_{\bl_\theta}$ such that
$$
\widehat\theta|_U=\bl_\theta,
\qquad
\widehat\theta|_{T^L_{2n-1}}=\theta.
$$
\end{enumerate}
\end{claim}
\begin{lemma}
Let
$$
\alpha_{2n-1}^E(\theta)=\ind_{\St_{\bl_\theta}}^{G(O)_{2n-1}}(\widehat\theta),
$$
and inflate this representation to $G(O)$. Then
\begin{enumerate}
\item $\alpha_{2n-1}^E(\theta)$ is irreducible, elliptic, special, and of depth $2n-1$;
\item $\alpha_{2n-1}^E(\theta)\simeq\alpha_{2n-1}^E(\theta^{-1})$.
\end{enumerate}
\end{lemma}
\begin{proof}
Irreducibility follows from Lemma \ref{MM}(2), and specialness follows from Theorem \ref{charodd}. The depth and ellipticity follow from the definition of $\bl_\theta$. The final assertion follows by conjugating with the nontrivial element of the normalizer of $T^L$.
\end{proof}
Hence the construction gives a map
$$
\bar\alpha_{2n-1}^E:\bar\Xi^L(2n-1)\longrightarrow\widehat G_{2n-1}^E.
$$
\begin{theorem}\label{od}
The map $\bar\alpha_{2n-1}^E$ is a bijection.
\end{theorem}
\begin{proof}
Surjectivity follows from Theorem \ref{charodd} and Lemma \ref{MM}(3): an elliptic representation recovers a character of its stabilizer, whose restriction to $T^L$ is the required parameter. If two parameters give isomorphic induced representations, their $U$-supports agree. Claim \ref{ellch}(3) then shows that the parameters differ only by inversion. Thus the map is injective.
\end{proof}
\subsection{Elliptic representations of even depth}\label{GO-even-depth}
We now treat positive even depth. Fix $n>0$. We denote by $T^L_m$ the image of $T^L$ in $G(O)_m$
(or in $G'(O)_m$, as appropriate), and by $T^L_m(r)$
the image of $T^L(r)$.
\begin{definition}\label{even-depth-data}
\begin{enumerate}
\item Put $A=G(O)_{2n}(n)$ and $U=G(O)_{2n}(n+1)$.
\item A character $\bl\in\widehat U$ is elliptic if its restriction to $G(O)_{2n}(2n)\simeq\fg(F)$ is nonzero and elliptic.
\item Let $\St_\bl$ be the stabilizer of $\bl$ in $G(O)_{2n}$.
\end{enumerate}
\end{definition}
\begin{lemma}\label{centralizer-lifting-even}
Let $R_r=F[t]/(t^r)$ and let $Y\in\fg(R_r)$ have regular semisimple constant term $Y_0$. If $T_0=Z_G(Y_0)$, then, after conjugation by an element of $\ker(G(R_r)\to G(F))$, one has
$$
Z_{G(R_r)}(Y)=T_0(R_r).
$$
\end{lemma}
\begin{proof}
Suppose the assertion has been obtained modulo $t^j$, and write the coefficient of $t^j$ of $Y$ as $Y_j$. Conjugation by $1+t^jX$ changes this coefficient by $[X,Y_0]$. Since $Y_0$ is regular semisimple and the trace form is nondegenerate,
$$
\fg(F)=\mathfrak z_\fg(Y_0)\oplus[\fg(F),Y_0].
$$
Thus $X$ can be chosen, uniquely modulo $\mathfrak z_\fg(Y_0)$, so that the new coefficient belongs to $\mathfrak z_\fg(Y_0)$. Induction on $j$ conjugates $Y$ into $\mathfrak z_\fg(Y_0)(R_r)$. If $g$ centralizes the resulting element, its reduction $g_0$ belongs to $T_0(F)$. Multiplying by $g_0^{-1}$ and comparing successive coefficients shows that every infinitesimal coefficient of $g$ belongs to $\mathfrak z_\fg(Y_0)$. Hence $g\in T_0(R_r)$.
\end{proof}
\begin{claim}\label{q}
Let $\bl\in\widehat U$ be elliptic.
\begin{enumerate}
\item $Z(A)=U$.
\item $[A,A]=G(O)_{2n}(2n)$.
\item There is an embedding $T^{L_\bl}\ho G(O)$ such that
$$
T^{L_\bl}_{2n}\subset\St_\bl,
\qquad
\St_\bl=T^{L_\bl}_{2n}A.
$$
\item One has
$$
[\St_\bl,\St_\bl]\cap T^{L_\bl}_{2n}=T^{L_\bl}_{2n}(2n).
$$
\item One has
$$
P(O)_{2n}\St_\bl=G(O)_{2n}.
$$
\end{enumerate}
\end{claim}
\begin{proof}
Modulo $t^{2n+1}$, the Lie algebra
$$
t^n\fg[[t]]/t^{2n+1}\fg[[t]]
$$
is two-step nilpotent because
$$
[t^n\fg,[t^n\fg,t^n\fg]]\subset t^{3n}\fg[[t]]\subset t^{2n+1}\fg[[t]].
$$
Since $2$ is invertible in $F$, the truncated exponential
$$
\exp_2(X)=1+X+\frac12X^2
$$
identifies $A$ with this two-step nilpotent Lie algebra, with multiplication
$$
X*Y=X+Y+\frac12[X,Y].
$$
The subspace $t^{n+1}\fg[[t]]/t^{2n+1}\fg[[t]]$ is central. Conversely, if the degree-$n$ coefficient of a central element were nonzero, its commutator with every element of $\fg(F)$ would vanish, contradicting $Z(\fg)=0$. This proves (1). Moreover,
$$
[\exp_2(t^nX),\exp_2(t^nY)]=\exp_2(t^{2n}[X,Y]),
$$
and $[\fg,\fg]=\fg$, proving (2).
Use the trace pairing to identify $\widehat U$ with $\fg(F[t]/(t^n))$, and let $Y$ correspond to $\bl$. An element of $G(O)_{2n}$ stabilizes $\bl$ if and only if its reduction modulo $t^n$ centralizes $Y$. Lemma \ref{centralizer-lifting-even} identifies this centralizer, after conjugating the torus by an element of $G(O)(1)$, with $T^{L_\bl}_n$. Since the kernel of the reduction map to this quotient is $A$, we obtain
$$
\St_\bl=T^{L_\bl}_{2n}A,
$$
which proves (3).
We already know that $[A,A]=G(O)_{2n}(2n)$. If $s\in T^{L_\bl}_{2n}$ and $a\in A$, the toral component of $[s,a]$ vanishes: on the Lie algebra this follows from
$$
(\Ad(s)-1)\fg\subset(\tori_\bl)^\perp.
$$
It follows that
$$
[T^{L_\bl}_{2n}A,T^{L_\bl}_{2n}A]\cap T^{L_\bl}_{2n}
=T^{L_\bl}_{2n}(2n),
$$
proving (4).
For (5), reduction modulo $t$ reduces the assertion to $P(F)T^{E_\bl}=G(F)$. The torus $T^{E_\bl}=E_\bl^*/F^*$ acts simply transitively on $G(F)/P(F)=\mathbb P^1(F)$: after choosing an $F$-basis of $E_\bl$, multiplication by a nonzero element sends the line $F\cdot1$ to an arbitrary $F$-line, and its stabilizer is $F^*$. Hence the equality holds modulo $t$. Suppose it has been lifted modulo $t^j$. The obstruction to lifting it modulo $t^{j+1}$ belongs to
$$
\fg/(\mathfrak p+\tori_\bl),
$$
which is zero because $\mathfrak p\cap\tori_\bl=0$ and the dimensions are $2$ and $1$. Induction gives $P(O)_{2n}T^{L_\bl}_{2n}=G(O)_{2n}$ and hence (5).
\end{proof}
Since $U$ is an $F$-vector group, there is a unique $F$-linear functional
$$
\ell_\bl:U\longrightarrow F
$$
such that
$$
\bl(u)=\widetilde\psi(\ell_\bl(u)).
$$
Put
$$
K_\bl=\ker(\ell_\bl),
\qquad
A_\bl=A/K_\bl.
$$
The quotient $A_\bl$, rather than $A/\ker(\bl)$, has a canonical one-dimensional additive central quotient. Indeed, every element of $A$ has a unique truncated logarithm
$$
t^nX_n+t^{n+1}X_{n+1}+\cdots+t^{2n}X_{2n}.
$$
The coefficient $X_n$ gives the first coordinate in $\fg(F)$. All remaining coefficients lie in the central subgroup $U$, and modulo $K_\bl$ their only invariant is the scalar obtained by applying $\ell_\bl$. Thus, using the section $X\mapsto\exp_2(t^nX)$, we identify $A_\bl$ as a set with $\fg(F)\times F$. Since triple commutators have degree at least $3n\geq2n+1$, the truncated Baker--Campbell--Hausdorff formula gives the multiplication
$$
(X,a)(Y,b)=\left(X+Y,a+b+\frac12\ell_{\bl,2n}([X,Y])\right),
$$
where $\ell_{\bl,2n}$ is the restriction of $\ell_\bl$ to $G(O)_{2n}(2n)\simeq\fg(F)$.
Use the trace pairing to write
$$
\ell_{\bl,2n}(Z)=\operatorname{tr}(y_\bl Z)
$$
for a unique regular elliptic element $y_\bl\in\fg(F)$. Put
$$
\tori_\bl=\mathfrak z_\fg(y_\bl),
\qquad
W_\bl=\fg/\tori_\bl.
$$
\begin{lemma}\label{B-symplectic}
The formula
$$
B_\bl(\overline X,\overline Y)=\operatorname{tr}(y_\bl[X,Y])
$$
defines a nondegenerate alternating form on $W_\bl$.
\end{lemma}
\begin{proof}
If $T\in\tori_\bl$, invariance of the trace pairing gives
$$
\operatorname{tr}(y_\bl[T,Y])=\operatorname{tr}([y_\bl,T]Y)=0,
$$
so the formula is well defined. If $B_\bl(\overline X,\overline Y)=0$ for every $Y$, then
$$
\operatorname{tr}([y_\bl,X]Y)=0
$$
for every $Y$. The trace pairing is nondegenerate, hence $[y_\bl,X]=0$ and $X\in\tori_\bl$.
\end{proof}
Choose the trace-orthogonal decomposition
$$
\fg=\tori_\bl\oplus W,
\qquad W=\tori_\bl^\perp,
$$
and transport $B_\bl$ to $W$. In coordinates $(T,w,a)\in\tori_\bl\times W\times F$, multiplication becomes
$$
(T,w,a)(T',w',a')=
\left(T+T',w+w',a+a'+\frac12B_\bl(w,w')\right).
$$
Consequently,
$$
A_\bl\simeq(\tori_\bl,+)\times R_W,
$$
where $R_W=W\times F$ is the Heisenberg group associated with $B_\bl$. In particular,
$$
p_\bl:A_\bl\longrightarrow\tori_\bl,
\qquad
p_\bl(T,w,a)=T,
$$
is a group homomorphism and $\ker(p_\bl)=R_W$.
\subsubsection{The stabilizer and its quasi-Heisenberg structure}
Put
$$
H_\bl=\St_\bl/K_\bl.
$$
Then $A_\bl$ is a normal subgroup of $H_\bl$. Since $\St_\bl$ fixes $\ell_\bl$ and $y_\bl$, it preserves $\tori_\bl$, $W$, and $B_\bl$. Thus there is a homomorphism
$$
a_\bl:H_\bl\longrightarrow\Sp(W,B_\bl).
$$
The subgroup $A_\bl$ acts trivially on $W$, as does the positive-loop subgroup of $T^{L_\bl}_{2n}$. The image of $a_\bl$ is therefore the compact anisotropic torus $T^{E_\bl}$. A splitting of the metaplectic cover over this torus gives a lift
$$
\widetilde a_\bl:H_\bl\longrightarrow\widetilde{\Sp(W)}.
$$
Let $C_\bl\subset T^{L_\bl}_{2n}$ be the image of $T^{L_\bl}(1)$.
\begin{lemma}\label{positive-loop-central}
The subgroup $C_\bl$ is contained in $Z(H_\bl)$.
\end{lemma}
\begin{proof}
Let $c\in C_\bl$. Its conjugation action on $R_W$ induces the identity on both $W=R_W/Z(R_W)$ and $Z(R_W)$. Every such automorphism of $R_W$ is inner, so it is conjugation by $(w_c,0)$ for a unique $w_c\in W$. Since $c$ commutes with $T^{E_\bl}$, the vector $w_c$ is fixed by $T^{E_\bl}$. The $T^{E_\bl}$-module $W$ has no nonzero fixed vector because the centralizer of $T^{E_\bl}$ in $\fg$ is $\tori_\bl$. Hence $w_c=0$, and $c$ centralizes $R_W$. It also commutes with $T^{L_\bl}_{2n}$ and with $(\tori_\bl,+)$. These subgroups generate $H_\bl$, so $c$ is central.
\end{proof}
The constant-loop torus $T^{E_\bl}$ acts faithfully on $W$, whereas $Z(H_\bl)R_W$ acts trivially on $W$. Claim \ref{q}(3), Lemma \ref{positive-loop-central}, and the decomposition $A_\bl=(\tori_\bl,+)R_W$ give a bijection
$$
T^{E_\bl}\times Z(H_\bl)R_W\stackrel{\sim}{\longrightarrow}H_\bl.
$$
Let $\rho$ be an irreducible representation of $\St_\bl$ on which $U$ acts by $\bl$. It factors through $H_\bl$. Let $\zeta_\rho$ be its central character. Then
$$
(H_\bl,R_W,T^{E_\bl},\widetilde a_\bl,\zeta_\rho)
$$
is a quasi-Heisenberg group.
\begin{proposition}\label{A-prime}
The restriction $\rho|_{R_W}$ is irreducible.
\end{proposition}
\begin{proof}
By Stone--von Neumann and Claim \ref{hei}(2),
$$
\rho|_{R_W}\simeq M\otimes\tau_W,
\qquad
M=\Hom_{R_W}(V_W,\rho).
$$
The action of $H_\bl$ on $M$ is initially projective; its cocycle is the inverse of the metaplectic cocycle occurring in the action on $\tau_W$. The chosen lift $\widetilde a_\bl$ cancels this cocycle. Consequently $M$ is an ordinary smooth representation of
$$
H_\bl/Z(H_\bl)R_W\simeq T^{E_\bl}.
$$
Irreducibility of $\rho$ implies irreducibility of $M$. Since the quotient torus is abelian, $M$ is one-dimensional. Hence $\rho|_{R_W}\simeq\tau_W$ is irreducible.
\end{proof}
\subsubsection{The Borel graph subgroup}
Let $\mathfrak p\subset\fg$ be the Lie algebra of $P$.
\begin{lemma}\label{p-to-W}
The natural map
$$
\mathfrak p\longrightarrow\fg/\tori_\bl
$$
is an isomorphism.
\end{lemma}
\begin{proof}
A nonzero element of the elliptic line $\tori_\bl$ is elliptic semisimple, whereas every element of $\mathfrak p$ has its eigenvalues in $F$. Hence $\mathfrak p\cap\tori_\bl=0$. Both spaces have dimension $2$.
\end{proof}
Let $R_P\subset A_\bl$ be the image of $P(O)_{2n}(n)$. Under
$$
A_\bl\simeq(\tori_\bl,+)\times R_W,
$$
the subgroup $R_P$ has the form
$$
R_P=\{(c(w),w,a):w\in W,\ a\in F\}
$$
for an $F$-linear map $c:W\to\tori_\bl$. Its commutator pairing is $B_\bl$, so $R_P$ is a Heisenberg group with center $F$.
\begin{proposition}\label{RP-irred}
The restriction $\rho|_{R_P}$ is irreducible.
\end{proposition}
\begin{proof}
By Proposition \ref{A-prime}, $\rho|_{R_W}=\tau_W$. The central group $(\tori_\bl,+)$ acts through a character $\eta$. Hence
$$
\rho(c(w),w,a)=\eta(c(w))\tau_W(w,a).
$$
There is an $F$-linear functional $d:W\to F$ such that $\eta(c(w))=\widetilde\psi(d(w))$. By nondegeneracy of $B_\bl$, there is $w_0\in W$ with $d(w)=B_\bl(w_0,w)$. Conjugation by $\tau_W(w_0,0)$, with the sign of $w_0$ adjusted if necessary, removes this twist. Thus $\rho|_{R_P}$ is isomorphic to the Heisenberg representation and is irreducible.
\end{proof}
\subsubsection{The finite-level restriction map}
Put $R_n=F[t]/(t^n)$. The trace pairing identifies $\widehat U$ with $\fg(R_n)$. If
$$
Y(t)=\sum_{r=0}^{n-1}Y_rt^r
$$
and
$$
X(t)=\sum_{j=n+1}^{2n}X_jt^j,
$$
then the corresponding character is
$$
\bl_Y(1+X(t))=
\widetilde\psi\left(\operatorname{coeff}_{t^{2n}}\operatorname{tr}(Y(t)X(t))\right).
$$
Under this identification, the coadjoint action is the adjoint action on $\fg(R_n)$, and ellipticity means that $Y_0$ is elliptic.
\begin{lemma}\label{orbit-restriction-injective}
Let $\mcO\subset\fg(R_n)$ be an adjoint orbit whose constant term is elliptic. Then restriction gives an injective map
$$
\mcO\longrightarrow\mathfrak p(R_n)^*.
$$
\end{lemma}
\begin{proof}
Write
$$
Y=\begin{pmatrix}a&b\\c&-a\end{pmatrix},
\qquad
Y'=\begin{pmatrix}a'&b'\\c'&-a'\end{pmatrix}
$$
with entries in $R_n$. Equality of their restrictions to $\mathfrak p(R_n)$ gives $a'=a$ and $c'=c$. Since $Y$ and $Y'$ are conjugate, their invariants
$$
q(Y)=-\det(Y)=a^2+bc
$$
and $q(Y')$ agree. The constant term $c_0$ is nonzero: otherwise $q(Y_0)=a_0^2$ would be a square, contrary to ellipticity. Thus $c$ is a unit in $R_n$, and
$$
b=\frac{q(Y)-a^2}{c}.
$$
The same formula gives $b'=b$.
\end{proof}
\begin{corollary}\label{finite-stabilizer-identity}
For elliptic $\bl\in\widehat U$,
$$
\St_{P(O)_{2n}}\left(\bl|_{U\cap P(O)_{2n}}\right)
=P(O)_{2n}\cap\St_{G(O)_{2n}}(\bl).
$$
\end{corollary}
\begin{proof}
Only the inclusion from left to right needs proof. If $p\in P(O)_{2n}$ stabilizes the restriction of $\bl$, then $p\bl$ and $\bl$ lie in the same $G(O)_{2n}$-orbit and have the same restriction. Lemma \ref{orbit-restriction-injective} gives $p\bl=\bl$.
\end{proof}
\begin{theorem}\label{charev}
Every elliptic representation of $G(O)$ of positive even depth is special.
\end{theorem}
\begin{proof}
Let $\pi$ be elliptic of depth $2n$. By Lemma \ref{MM}, there are an elliptic $\bl\in\widehat U$ and an irreducible representation $\rho$ of $\St_\bl$ such that $U$ acts on $\rho$ by $\bl$ and
$$
\pi\simeq\ind_{\St_\bl}^{G(O)_{2n}}(\rho),
$$
followed by inflation to $G(O)$. Put $H'=P(O)_{2n}$. By Proposition \ref{RP-irred}, $\rho|_{R_P}$ is irreducible. Since $R_P\subset H'\cap\St_\bl$, it follows that $\rho|_{H'\cap\St_\bl}$ is irreducible. Claim \ref{q}(5) gives $H'\St_\bl=G(O)_{2n}$, and Lemma \ref{orbit-restriction-injective} supplies the injectivity hypothesis of Lemma \ref{3.12}. Therefore $\pi|_{P(O)}$ is irreducible.
\end{proof}
\subsubsection{Parametrization at even depth}
Let $E/F$ be quadratic, put $L=E((t))$, and let $\theta\in\Xi^L(2n)$. Since $2n<2(n+1)$, Claim \ref{iota} applies to $p_{2n}(n+1)$. Define
$$
\bl_\theta=\theta\circ p_{2n}(n+1)\in\widehat U.
$$
Then $\bl_\theta$ is elliptic and has depth $2n$. Let
$$
H_\theta=\St_{\bl_\theta}/K_{\bl_\theta},
$$
and retain the notation $R_W$, $C_{\bl_\theta}$, and $T^E$ from the preceding construction.
\begin{lemma}\label{central-character-even}
One has
$$
Z(H_\theta)=C_{\bl_\theta}Z(R_W).
$$
The character $\theta$ on $C_{\bl_\theta}$ and the central character $\bl_\theta$ on $Z(R_W)$ agree on their intersection and therefore define a unique character
$$
\zeta_\theta:Z(H_\theta)\longrightarrow\CC^*.
$$
\end{lemma}
\begin{proof}
The radical $(\tori_{\bl_\theta},+)$ of the commutator pairing is the image of the appropriate congruence subgroup of $T^L(1)$, hence is contained in $C_{\bl_\theta}$. Lemma \ref{positive-loop-central} then gives the displayed description of the center. On the intersection, the map $p_{2n}(n+1)$ is the identity on the torus, so the two characters agree.
\end{proof}
Choose a splitting $s:T^E\to\widetilde{\Sp(W)}$. It gives a lift $\widetilde a_{\bl_\theta,s}:H_\theta\to\widetilde{\Sp(W)}$. Applying Corollary \ref{alpha} with central character $\zeta_\theta$ and with the character $\theta|_{T^E}$ of the quotient torus gives an irreducible representation $\rho_{\theta,s}$ of $\St_{\bl_\theta}$. Put
$$
\alpha_{2n,s}^E(\theta)
=
\ind_{\St_{\bl_\theta}}^{G(O)_{2n}}(\rho_{\theta,s}),
$$
and inflate to $G(O)$.
\begin{proposition}\label{even-fixed-splitting}
For every splitting $s$, the construction has the following properties.
\begin{enumerate}
\item The representation $\alpha_{2n,s}^E(\theta)$ is irreducible, elliptic, special, and of depth $2n$.
\item One has
$$
\alpha_{2n,s}^E(\theta)\simeq\alpha_{2n,s}^E(\theta^{-1}).
$$
\item The induced map
$$
\bar\alpha_{2n,s}^E:\bar\Xi^L(2n)\longrightarrow\widehat G_{2n}^E
$$
is a bijection.
\end{enumerate}
\end{proposition}
\begin{proof}
Irreducibility follows from Lemma \ref{MM}(2), and specialness from Theorem \ref{charev}. The depth and ellipticity are read from $\bl_\theta$. Galois conjugation is implemented by the nontrivial normalizer element and sends $\theta$ to $\theta^{-1}$, proving (2).
Conversely, let $\pi\in\widehat G_{2n}^E$. Mackey theory gives an elliptic $\bl$ and an irreducible representation $\rho$ of $\St_\bl$ such that
$$
\pi\simeq\ind_{\St_\bl}^{G(O)_{2n}}\rho.
$$
By Claim \ref{q}, after conjugacy we may identify the torus in the stabilizer with $T^L_{2n}$. The action of the positive-loop subgroup is central in $H_\bl$ and hence defines a character $\theta_1$ of $T^L(1)$. Proposition \ref{A-prime} and the fixed splitting $s$ identify the Weil multiplicity space with a character $\theta_0$ of $T^E$. Their restrictions to the intersection agree, since both equal the central character of $\rho$. Thus $\theta_0$ and $\theta_1$ glue to a character $\theta$ of
$$
T^L=T^E T^L(1).
$$
Its restriction to the last nontrivial congruence quotient is $\bl$, so $\theta\in\Xi^L(2n)$. This construction is inverse to $\theta\mapsto\alpha^E_{2n,s}(\theta)$: the central character recovers $\theta_1$, while the Weil multiplicity space recovers $\theta_0$.
Finally, Mackey's isomorphism criterion shows that two induced representations are isomorphic precisely when their inducing data are conjugate. The quotient of the normalizer of $T^L$ by $T^L$ has order two, and its nontrivial element acts on $T^L=L^*/K^*$ by Galois conjugation, which is inversion. Hence the only identification is $\theta\sim\theta^{-1}$. This proves (3).
\end{proof}
\subsubsection{The unramified case}
If $E/F$ is unramified, Corollary \ref{lift} gives a canonical splitting over $T^E\simeq E^1$. Proposition \ref{even-fixed-splitting} therefore gives a canonical bijection
$$
\bar\alpha_{2n}^E:\bar\Xi^L(2n)\stackrel{\sim}{\longrightarrow}\widehat G_{2n}^E.
$$
\subsubsection{The ramified case}
Assume that $E/F$ is ramified. By Lemma \ref{tori}, the set of splittings over $T^E$ is a torsor under
$$
\D_E=\Hom(T^E,\{\pm1\})\simeq\ZZ/2\ZZ.
$$
Let $\delta$ be its nontrivial element. If $s$ is a splitting, the other splitting is $\delta s$, and Corollary \ref{alpha} gives
$$
\alpha_{2n,\delta s}^E(\theta)
\simeq
\alpha_{2n,s}^E(\delta\theta),
$$
where $\delta$ is extended trivially across $T^L(1)$. Thus the class of the parameter in
$$
\ti\Xi^L(2n)=\bar\Xi^L(2n)/\D_L
$$
is independent of the splitting.
The action of $\delta$ on $\bar\Xi^L(2n)$ is free. Indeed, if $[\delta\theta]=[\theta]$, then either $\delta\theta=\theta$, which gives $\delta=1$, or $\delta\theta=\theta^{-1}$, which gives $\theta^2=\delta$. In the second case $\theta|_{T^L(1)}$ has order at most two, because $\delta$ is trivial on $T^L(1)$. But $T^L(1)$ is pro-$p$ and $p\neq2$, so it has no nontrivial character of order two. Hence $\theta$ is trivial on $T^L(1)$, contradicting the positive-depth assumption.
The involution on $\widehat G_{2n}^E$ may be described as follows. Choose $s$, write $\pi=\bar\alpha_{2n,s}^E([\theta])$, and put
$$
\delta\cdot\pi=\bar\alpha_{2n,s}^E([\delta\theta]).
$$
The preceding transformation formula shows that this definition is independent of $s$. It is free, and its orbits are exactly the fibers of the canonical map
$$
\widehat G_{2n}^E\longrightarrow\ti\Xi^L(2n).
$$
Hence $\widehat G_{2n}^E$ is naturally a $\D_L$-torsor over $\ti\Xi^L(2n)$.
\begin{theorem}\label{equnr}
Let $E/F$ be quadratic and put $L=E((t))$.
\begin{enumerate}
\item At depth $0$, the classical classification gives a natural bijection
$$
\bar\Xi^L(0)\stackrel{\sim}{\longrightarrow}\widehat G_0^E.
$$
\item If $m>0$ is odd, there is a natural bijection
$$
\bar\Xi^L(m)\stackrel{\sim}{\longrightarrow}\widehat G_m^E.
$$
\item If $m>0$ is even and $E/F$ is unramified, there is a natural bijection
$$
\bar\Xi^L(m)\stackrel{\sim}{\longrightarrow}\widehat G_m^E.
$$
\item If $m>0$ is even and $E/F$ is ramified, there is a canonical map
$$
\widehat G_m^E\longrightarrow\ti\Xi^L(m)
$$
whose fibers are torsors under $\D_L\simeq\ZZ/2\ZZ$.
\end{enumerate}
\end{theorem}

\begin{corollary}\label{cunr}
Let $\pi$ be a special representation of $G(O)$ of positive depth $m$. Then there are a closed almost unipotent subgroup $R\subset G(O)_m$ and an irreducible representation $\rho$ of $R$ such that $\pi$ is the inflation of
$$
\ind_R^{G(O)_m}(\rho).
$$
At odd depth, $\rho$ is a character. At even depth, $\rho$ is the quasi-Heisenberg representation described above.
\end{corollary}
\begin{proof}
At odd depth take $R=\St_{\bl_\pi}$ and $\rho=\rho_\pi$, which is a character. At even depth take the stabilizer and its irreducible quasi-Heisenberg representation. These stabilizers are extensions of anisotropic tori by unipotent groups and are therefore almost unipotent.
\end{proof}
\section{Special representations of $G'(O)$}\label{G'O}
We now carry out the analogous analysis for $G'(O)$. Let
$$
\kk:G'(O)\longrightarrow\ZZ/2\ZZ\ltimes F^*
$$
be the homomorphism for which the nontrivial element of $\ZZ/2\ZZ$ acts on $F^*$ by inversion, and which is defined as follows. If $g\in I$ is represented by $\begin{pmatrix}a&b\\c&d\end{pmatrix}$ and $a_0,d_0$ are the constant terms of $a,d$, put
$$
\kk(g)=(1,a_0/d_0).
$$
For $s=\begin{pmatrix}0&1\\t&0\end{pmatrix}$, put $\kk(s)=(-1,1)$.
\begin{claim}\label{act}
\begin{enumerate}
\item The subgroups $G'(O)(m)$ are normal in $G'(O)$.
\item The subquotients $G'(O)_m(m)$ are naturally $F$-vector spaces.
\item If $m\in\ZZ$, then $G'(O)_m(m)$ is one-dimensional. The element $s$ acts by $-1$, and $P(O)$ acts trivially.
\item If $m\in\frac12\ZZ\sm\ZZ$, then
$$
G'(O)_m(m)=M_m^+\oplus M_m^-,
$$
where
$$
M_m^+\simeq\fm^{m-1/2}/\fm^{m+1/2},
\qquad
M_m^-\simeq\fm^{m+1/2}/\fm^{m+3/2}.
$$
Each summand is one-dimensional. The choice of $t$ gives generators $x_m^+$ and $x_m^-$ and identifies the subquotient with $F\oplus F$.
\item The adjoint action on $M_m^+\oplus M_m^-$ factors through $\kk$: an element $a\in F^*$ acts by $\operatorname{diag}(a,a^{-1})$, and $s$ interchanges the two summands.
\end{enumerate}
\end{claim}
\begin{definition}\label{ac}
Let $m=n-1/2$ with $n\in\ZZ_{>0}$.
\begin{enumerate}
\item If $\mu$ is a character of $G'(O)_m(m)=M_m^+\oplus M_m^-$, define $[\mu^\pm]\in F$ by
$$
\mu(x_m^\pm z)=\widetilde\psi([\mu^\pm]z),\qquad z\in F.
$$
\item A character $\mu$ is elliptic if it restriction on $M_m^+$ and on $M_m^-$ does not vanish.
\item Put
$$
\det(\mu)=t[\mu^+][\mu^-]\in K.
$$
For an elliptic character, this is nonzero; its square class is independent of the chosen uniformizer and generators.
\item For an irreducible representation $\pi$ of $G'(O)$, let $m(\pi)$ be the least $m\in\frac12\ZZ_{\geq0}$ such that $G'(O)(m+1/2)$ acts trivially.
\item The representation $\pi$ is special if it has dimension greater than $1$ and $\pi|_{P(O)}$ is irreducible.
\item A non-one-dimensional irreducible representation $\pi$ is elliptic if $m(\pi)\in\frac12\ZZ\sm\ZZ$ and the support of its restriction to each of $M_m^+$ and $M_m^-$ is nontrivial.
\end{enumerate}
\end{definition}


\begin{remark}
\begin{enumerate}
\item Replacing $t$ by another uniformizer and changing the generators of $M_m^\pm$ changes $t[\mu^+][\mu^-]$ by a square. Thus the quadratic extension $K(\sqrt{t[\mu^+][\mu^-]})$ is intrinsic.
\item Elliptic representations of $G'(O)$ necessarily have positive half-integral depth.
\end{enumerate}
\end{remark}
\begin{lemma}\label{Gprime-nonelliptic}
Let $\pi$ be an irreducible representation of $G'(O)$.
\begin{enumerate}
\item If $m(\pi)\in\ZZ$, then $\pi|_{P(O)}$ is reducible.
\item If $m(\pi)\notin\ZZ$ and $\pi$ is not elliptic, then $\pi|_{P(O)}$ is reducible.
\end{enumerate}
\end{lemma}
\begin{proof}
The case $m=0$ is easy, so we assume that $m>0$.
Suppose first that $m=m(\pi)\in\ZZ$. The support of $\pi$ on the nonzero one-dimensional group $G'(O)_m(m)$ is a nontrivial $G'(O)$-orbit. By Claim \ref{act}, $P(O)$ acts trivially on this character group, whereas $s$ acts by sign. Thus the support contains at least two $P(O)$-orbits, and Lemma \ref{MM}(1) implies reducibility.
Now let $m=n-1/2$. If $\pi$ is not elliptic, its nonzero support on $M_m^+\oplus M_m^-$ is the union
$$
((M_m^+)^*\sm0)\sqcup((M_m^-)^*\sm0).
$$
The group $G'(O)$ interchanges the two axes, but $P(O)$ preserves each. Hence the support splits into two $P(O)$-orbits, and the restriction is reducible by Lemma \ref{MM}(1).
\end{proof}
\begin{corollary}
Every special representation of $G'(O)$ is elliptic.
\end{corollary}
\subsection{Elliptic representations of $G'(O)$}
Fix $n\in\ZZ_{>0}$, put $m=n-1/2$, and let $U$ be the image of $G'(O)(n/2)$ in $G'(O)_m$.
\begin{claim}\label{Gprime-U}
The group $U$ is a closed normal commutative subgroup of $G'(O)_m$.
\end{claim}
\begin{definition}\label{c}
Let $\bl\in\widehat U$.
\begin{enumerate}
\item Let $\St_\bl$ be its stabilizer in $G'(O)_m$.
\item Let $\bar\bl$ be its restriction to $G'(O)_m(m)=M_m^+\oplus M_m^-$.
\item The character $\bl$ is elliptic if $[\bar\bl^+]$ and $[\bar\bl^-]$ are both nonzero.
\item For elliptic $\bl$, put
$$
L_\bl=K\left(\sqrt{t[\bar\bl^+][\bar\bl^-]}\right).
$$
\end{enumerate}
\end{definition}
\begin{claim}\label{ch}
Let $\bl$ be elliptic and put $L=L_\bl$.
\begin{enumerate}
\item For a suitable embedding $T^L\ho G'(O)$, one has
$$
\St_\bl=T^L_mU.
$$
\item One has
$$
P(O)_m\St_\bl=G'(O)_m.
$$
\item The restriction map from the $G'(O)_m$-orbit of $\bl$ to $\widehat{U\cap P(O)_m}$ is injective.
\end{enumerate}
\end{claim}
\begin{proof}
The leading pair $([\bar\bl^+],[\bar\bl^-])$ has both coordinates nonzero. Claim \ref{act}(5) shows that its stabilizer is the anisotropic torus attached to the square class of $t[\bar\bl^+][\bar\bl^-]$. Successive approximation through the filtration gives (1). The same coordinate calculation shows that $P(O)_m$ is transitive on the orbit and that restriction to $U\cap P(O)_m$ determines both coordinates, giving (2) and (3).
\end{proof}
\begin{theorem}\label{char}
Every elliptic representation of $G'(O)$ is special.
\end{theorem}
\begin{proof}
Let $\pi$ be elliptic of depth $m=n-1/2$. By Lemma \ref{MM}, there are an elliptic character $\bl$ of $U$ and an irreducible representation $\rho$ of $\St_\bl$ such that $U$ acts on $\rho$ by $\bl$ and
$$
\pi\simeq\ind_{\St_\bl}^{G'(O)_m}(\rho),
$$
followed by inflation to $G'(O)$. By Claim \ref{ch}(1), the quotient $\St_\bl/\ker(\bl)$ is commutative, so $\rho$ is a character. Claim \ref{ch}(2)--(3) and Lemma \ref{3.12} imply that $\pi|_{P(O)}$ is irreducible.
\end{proof}
\subsubsection{Classification}
For a ramified quadratic extension $L/K$, let $\widehat G_m^L$ be the set of isomorphism classes of elliptic representations of $G'(O)$ of depth $m$ with associated extension $L$.
An elliptic representation $\pi$ determines, through the preceding proof, a character of $T^L_m$ and hence a character $\theta_\pi$ of $T^L$. Changing the embedding of $T^L$ by the nontrivial normalizer element replaces $\theta_\pi$ by its Galois conjugate, which is its inverse. Thus there is a recovery map
$$
\beta_m^L:\widehat G_m^L\longrightarrow\bar\Xi^L(m).
$$
\begin{claim}\label{Gprime-parameter}
Conversely, for every $\theta\in\Xi^L(m)$ there are an elliptic character $\bl_\theta\in\widehat U$ and a unique character $\widehat\theta$ of $\St_{\bl_\theta}$ such that
$$
L_{\bl_\theta}=L,
\qquad
\widehat\theta|_U=\bl_\theta,
\qquad
\widehat\theta|_{T^L_m}=\theta.
$$
The construction is compatible with inversion.
\end{claim}
Define
$$
\alpha_m^L(\theta)=\ind_{\St_{\bl_\theta}}^{G'(O)_m}(\widehat\theta),
$$
followed by inflation to $G'(O)$.
\begin{proposition}\label{Oodd}
The construction induces a bijection
$$
\bar\alpha_m^L:\bar\Xi^L(m)\stackrel{\sim}{\longrightarrow}\widehat G_m^L,
$$
whose inverse is $\beta_m^L$.
\end{proposition}
\begin{proof}
Lemma \ref{MM}(2) gives irreducibility, and Theorem \ref{char} gives specialness. The support recovers $\bl_\theta$ up to the normalizer action, and the stabilizer character then recovers $\theta$ up to inversion. Thus the construction and recovery maps are inverse.
\end{proof}
\begin{theorem}\label{cunr-summary}
\begin{enumerate}
\item Let $\pi$ be a special representation of $G(O)$ of depth $m$. If $m>0$ then $\pi$ is the inflation of $\ind_R^{G(O)_m}(\rho)$ for a closed almost unipotent subgroup $R\subset G(O)_m$ and an irreducible representation $\rho$ of $R$. At odd depth $\rho$ may be taken to be a character; at positive even depth it is the quasi-Heisenberg representation of Section \ref{GO-even-depth}.
\item One has $\pi|_{P(O)}\simeq\bs(m)$.
\item Let $\pi$ be a special representation of $G'(O)$ of depth $n-1/2$. Then $\pi$ is the inflation of $\ind_R^{G'(O)_{n-1/2}}(\nu)$ for a closed almost unipotent subgroup $R$ and a character $\nu$ of $R$.
\item One has $\pi|_{P(O)}\simeq\bs(n-1)$.
\end{enumerate}
\end{theorem}
\begin{remark}The first two statements restate Corollary 5.33 and Lemma 4.7.\end{remark}
%------------------------------------------------------------------------------------------------------------------
\section{Some results on abelian categories}
In this section we formulate and prove some results from category theory needed later.
\sms
Let $\mcC$ be an abelian category.
\begin{definition}\label{strict}
An object $\mV$ of $\Pro(\mcC)$ is strict if it admits a presentation
$$
\mV\simeq\prl_{n\geq0}V_n,\qquad V_n\in\operatorname{Ob}(\mcC),
$$
whose transition morphisms $V_{n+1}\to V_n$ are epimorphisms.
\end{definition}
\begin{remark}
Since $\mcC$ has finite limits, $\Pro(\mcC)$ has all small limits. In particular, the products used below exist.
\end{remark}
\begin{definition}\label{basic}
Let $(\mcA,\times,*)$ be a category with finite products. A contravariant action of $\mcA$ on a category $\mcC$ consists of the following data:
\begin{enumerate}
\item A functor $a:\mcA^{op}\times\mcC\to\mcC$, together with a functorial identification $a(*,V)\simeq V$.
\item A functorial isomorphism
$$
\alpha_{X,Y,V}:a(X\times Y,V)\xrightarrow{\sim}a(X,a(Y,V)),
$$
for $X,Y\in\mcA$ and $V\in\mcC$, satisfying the usual associativity and unit axioms.
\end{enumerate}
\end{definition}
\begin{definition}
\begin{enumerate}
\item A group object in $\mcA$ is an object $G$ equipped with morphisms
$$
m:G\times G\to G,\qquad e:*\to G,\qquad \iota:G\to G,
$$
satisfying the usual associativity, unit, and inverse axioms.
\item Let $G$ be a group object of $\mcA$ and let $V\in\operatorname{Ob}(\mcC)$. A representation of $G$ on $V$ is a morphism
$$
\pi:V\to a(G,V)
$$
such that
$$
\alpha_{G,G,V}^{-1}\circ a(G,\pi)\circ\pi=a(m,V)\circ\pi
$$
and
$$
a(e,V)\circ\pi=\operatorname{id}_V,
$$
where we use the identification $a(*,V)\simeq V$.
\item A representation of $G$ on $V$ is irreducible if $V$ has no nonzero proper $G$-invariant subobject.
\item Similarly, an action of an abstract group $\GG$ on $V$ is irreducible if $V$ has no nonzero proper $\GG$-invariant subobject.
\end{enumerate}
\end{definition}
\subsection{Filtrations}
Let $\mV$ be a strict object of $\Pro(\mcC)$ and let $\mcF=\{\mW_i\}_{i\in\ZZ}$ be an increasing filtration of $\mV$ by subobjects.
\begin{definition}\label{admissible}
\begin{enumerate}
\item A subobject $\mU\subset\mV$ is dense with respect to $\mcF$ if, for every $i\in\ZZ$, the natural morphism
$$
\beta_i:\frac{\mW_{i+1}\cap\mU}{\mW_i\cap\mU}\longrightarrow\frac{\mW_{i+1}}{\mW_i}
$$
is an isomorphism.
\item The filtration $\mcF$ is admissible if, for every epimorphism $\kk:\mV\to V$ with $V\in\mcC$, there exist $r_-(\kk),r_+(\kk)\in\ZZ$ such that
$$
\kk(\mW_{r_-(\kk)})=0,\qquad \kk(\mW_{r_+(\kk)})=V.
$$
\item We put
$$
\overline{\mV}_{\mcF}:=\prod_{i\in\ZZ}\mW_{i+1}/\mW_i.
$$
\end{enumerate}
\end{definition}
\begin{claim}\label{separated}
If $\mcF$ is admissible, then
$$
\bigcap_{i\in\ZZ}\mW_i=0.
$$
\end{claim}
\begin{proof}
Choose a strict presentation $\mV\simeq\varprojlim_{n\geq0}V_n$ with epimorphic transition morphisms, and denote the canonical epimorphisms by $\kk_n:\mV\to V_n$. Put $\mI:=\bigcap_{i\in\ZZ}\mW_i$. For every $n$, admissibility gives an integer $r_-(n)$ such that $\kk_n(\mW_{r_-(n)})=0$. Since $\mI\subset\mW_{r_-(n)}$, the composite $\mI\to\mV\xrightarrow{\kk_n}V_n$ is zero. The morphisms $\kk_n$ are jointly monomorphic, and therefore $\mI=0$.
\end{proof}
\begin{lemma}\label{dense}
Let $\mcF=\{\mW_i\}_{i\in\ZZ}$ be an admissible filtration of a strict object $\mV$ of $\Pro(\mcC)$, and let $\mU\subset\mV$ be dense with respect to $\mcF$. Then $\mU=\mV$.
\end{lemma}
\begin{proof}
It is enough to show that, for every epimorphism $\kk:\mV\to V$ with $V\in\mcC$, one has $\kk(\mU)=V$. Put
$$
V_i:=\kk(\mW_i),\qquad U_i:=\kk(\mW_i\cap\mU).
$$
Choose $r_-=r_-(\kk)$ and $r_+=r_+(\kk)$ as in Definition \ref{admissible}. The isomorphisms $\beta_i$ imply that the natural morphisms
$$
U_{i+1}/U_i\longrightarrow V_{i+1}/V_i
$$
are surjective.
 Since $U_{r_-}=V_{r_-}=0$, induction on $i$ gives $U_i=V_i$ for $r_-\leq i\leq r_+$. Hence
$$
V=V_{r_+}=U_{r_+}\subset\kk(\mU),
$$
so $\kk(\mU)=V$. Applying this to the epimorphic projections of a strict presentation of $\mV$, we obtain $\mU=\mV$.
\end{proof}
\begin{lemma}\label{admi}
Let $\mV=\varprojlim_{l\geq0}V_l$, where $V_l\in\operatorname{Ob}(\mcC)$ and the transition morphisms are epimorphisms, and let $\kk_l:\mV\to V_l$ be the canonical projections. Let $\mcF=\{\mW_n\}_{n\in\ZZ}$ be an increasing filtration of $\mV$ such that, for every $l\geq0$, there exist integers $r_-(l)$ and $r_+(l)$ satisfying
$$
\kk_l(\mW_{r_-(l)})=0,\qquad \kk_l(\mW_{r_+(l)})=V_l.
$$
Then
\begin{enumerate}
\item The filtration $\mcF$ is admissible.
\item If $b:\mV'\to\mV$ is a morphism such that $\kk_l\circ b:\mV'\to V_l$ is an epimorphism for every $l\geq0$, then $b$ is an epimorphism.
\end{enumerate}
\end{lemma}
\begin{proof}
(1) Let $\kk:\mV\to V$ be an epimorphism with $V\in\operatorname{Ob}(\mcC)$. The morphism $\kk$ factors through some $\kk_l$, say $\kk=\tau\circ\kk_l$ for a morphism $\tau:V_l\to V$. Since $\kk$ is an epimorphism, so is $\tau$. Hence
$$
\kk(\mW_{r_-(l)})=0,\qquad \kk(\mW_{r_+(l)})=V,
$$
and $\mcF$ is admissible.

(2) Suppose that $f,g:\mV\to\mY$ satisfy $f\circ b=g\circ b$. Choose a presentation $\mY=\varprojlim_jY_j$ and let $p_j:\mY\to Y_j$ be the canonical projections. For every $j$, the morphism $p_j\circ(f-g):\mV\to Y_j$ factors through some $\kk_l$, say
$$
p_j\circ(f-g)=h\circ\kk_l.
$$
Since $(p_j\circ(f-g))\circ b=0$ and $\kk_l\circ b$ is an epimorphism, we have $h=0$. Thus $p_j\circ f=p_j\circ g$ for every $j$, and hence $f=g$. Therefore $b$ is an epimorphism.
\end{proof}




\subsection{Graded objects}
\begin{definition}
Let $\mcF=\{\mW_i\}_{i\in\ZZ}$ be a filtration of $\mV$, let $\GG$ be an abstract group, let
$$
\nu:\GG\to\ZZ,\qquad \Gg\mapsto[\Gg],
$$
be a homomorphism, and let $\pi:\GG\to\operatorname{Aut}(\mV)$ be a representation satisfying
$$
\pi(\Gg)(\mW_i)=\mW_{i+[\Gg]}.
$$
Then $\GG$ acts on
$$
\overline{\mV}_{\mcF}=\prod_{i\in\ZZ}\mW_{i+1}/\mW_i.
$$
We denote the resulting representation by $\overline\pi$.
\end{definition}
\begin{lemma}\label{bar}
Assume that $\mcF$ is admissible and that every nonzero $\GG$-invariant subobject $\mU\subset\mV$ satisfies
$$
\mU\cap\mW_i\neq0
$$
for some $i\in\ZZ$. If $\overline\pi$ is irreducible, then $\pi$ is irreducible.
\end{lemma}
\begin{proof}
Let $\mU\subset\mV$ be a nonzero $\GG$-invariant subobject and put
$$
\overline{\mU}_{\mcF}:=\prod_{i\in\ZZ}\frac{\mW_{i+1}\cap\mU}{\mW_i\cap\mU}.
$$
This is a $\GG$-invariant subobject of $\overline{\mV}_{\mcF}$. It is nonzero: otherwise all the subobjects $\mW_i\cap\mU$ would be equal, and Claim \ref{separated} would imply that they are all zero, contrary to the hypothesis on $\mU$. Since $\overline\pi$ is irreducible, we have
$$
\overline{\mU}_{\mcF}=\overline{\mV}_{\mcF}.
$$
Thus $\mU$ is dense with respect to $\mcF$, and Lemma \ref{dense} gives $\mU=\mV$.
\end{proof}

\sms

\section{A digression on representation theory of algebraic groups over $K$}\label{pro-vect} In this section we remind (slightly reformulated) definitions and constructions from \cite{GK}.
\subsection{Basic definitions}

\subsubsection{Categories}

Besides the large categories $Set$ of sets and $\Vect$ of complex vector spaces, we will only consider the following categories:

\begin{enumerate}

\item monoidal categories $\mcA$ obtained from the monoidal category $(Set_0,\otimes)$ of finite sets, where $\otimes$ is the Cartesian product, by the operations $\Ind_{\aleph_0}$ and $\Pro_{\aleph_0}$;

\item categories $\mcC$ obtained from the category $\Vect_0$ of finite-dimensional complex vector spaces by the operations $\Ind_{\aleph_0}$ and $\Pro_{\aleph_0}$.\footnote{In this paper we restrict ourselves to countable limits and colimits. This is mostly due to our laziness, but it is also sufficient for our purposes; for instance, all vector spaces which appear in this paper have countable dimension, and all pro-vector spaces are countable limits of such spaces.}

\end{enumerate}

More precisely, we consider the following lists of monoidal categories $\mcA$ and categories $\mcC$.

\begin{definition}\label{cat}

Categories $\mcA$:

\begin{enumerate}

\item $Set$ is the large category of sets.

\item $\bSet=\Ind_{\aleph_0}\Pro_{\aleph_0}(Set_0)$.

\item $\Pro\bSet:=\Pro_{\aleph_0}(\bSet)$.

\item $\mSet:=\Ind_{\aleph_0}(\Pro\bSet)$.

\end{enumerate}

Categories $\mcC$:

\begin{enumerate}

\item $\Vect$ is the large category of complex vector spaces.

\item $\Vect_{\aleph_0}=\Ind_{\aleph_0}(\Vect_0)\subset\Vect$ is the category of vector spaces of at most countable dimension.

\item $\mVect=\Pro_{\aleph_0}(\Vect_{\aleph_0})$.

\end{enumerate}

\end{definition}

\begin{remark}

\begin{enumerate}

\item $\bSet$ is a subcategory of the category of locally profinite topological spaces.

\item For an algebraic variety $X$ over $F$ we consider the set $X(F)$ as an object of $\bSet$.

\end{enumerate}

\end{remark}

\begin{proposition}\label{weak-strict-GK}
Let $H\in\Pro\bSet$ be a  strict group object. Then, in the countable setting used here, the natural functor
$$
\Pro_{\aleph_0}\bigl(\Rep(H,\Vect_{\aleph_0})\bigr)\longrightarrow\Rep(H,\mVect)
$$
is an equivalence.
\end{proposition}

\begin{proof}
This is the countable version of \cite[Proposition~2.5]{GK}.
\end{proof}

\begin{claim}

The category $\mVect$ has countable inductive and projective limits.

\end{claim}

\begin{definition}\label{maps}

\begin{enumerate}

\item For $\mV=\varprojlim V_n$, with $V_n\in\Vect_{\aleph_0}$, we write $\ti\mV:=\varprojlim V_n\in\Vect$, where $\varprojlim V_n$ is the projective limit in $\Vect$.

\item For $\mX\in\mSet$ we write $\ti\mX:=\Hom_{\mSet}(*,\mX)$.

\end{enumerate}

\end{definition}

\begin{claim}

\begin{enumerate}

\item The correspondence $\mV\mapsto\ti\mV$ defines a functor from $\mVect$ to $\Vect$, and the correspondence $\mX\mapsto\ti\mX$ defines a monoidal functor from $\mSet$ to $Set$.

\item If $\mT\in\mSet$ is a group object, then the set $\ti\mT$ has a group structure.

\item If $W$ is regarded as a discrete vector space, then there is a canonical isomorphism $\Hom_{\mVect}(\mV,W)\simeq\Hom_{\Vect}^{\mathrm{cont}}(\ti\mV,W)$, where $\ti\mV$ is endowed with its inverse-limit topology.

\end{enumerate}

\end{claim}

\begin{definition}

We define a functor $\un{\Hom}:\mVect\times\mVect\to\mVect$ by $\un{\Hom}(\mV,\mW)=\prl_m\inl_n\Hom(V_n,W_m)$ for $\mV=\prl_nV_n$ and $\mW=\prl_mW_m$.

\end{definition} \subsubsection{Actions}

\begin{definition}\label{a}

\begin{enumerate}

\item $a_0:Set_0^{op}\times\Vect\to\Vect$ is the functor such that $a_0(X,V)$ is the space of maps from $X$ to $V$. The functor $a_0$ is contravariant in $X$ and covariant in $V$.

\item $a_1:\bSet^{op}\times\Vect\to\Vect$ is the functor such that $a_1(\mathbf X,V)$ is the space of locally constant $V$-valued functions on $\mathbf X$.

\item $a_2:(\Pro\bSet)^{op}\times\Vect\to\Vect$ is the functor such that $a_2(\mX,V):=\inl_n a_1(\mathbf X_n,V)$, where $\mX=\prl_n\mathbf X_n$.

\item $a_3:(\Pro\bSet)^{op}\times\mVect\to\mVect$ is the functor such that $a_3(\mX,\mV):=\prl_n a_2(\mX,V_n)$, where $\mV=\prl_nV_n$.

\item $a:\mSet^{op}\times\mVect\to\mVect$ is the functor such that $a(\mX,\mV):=\prl_n a_3(\mX_n,\mV)$, where $\mX=\inl_n\mX_n$.

\end{enumerate}

\end{definition}

\begin{lemma}\label{action}

\begin{enumerate}

\item The functor $a$ defines an action of the category $\mSet$ on $\mVect$, in the contravariant convention of Definition \ref{basic}.

\item The functors $\mX\mapsto\ti\mX$ and $\mV\mapsto\ti\mV$ are compatible with the action $a$ in the sense that an action of a group object $\mT$ on $\mV$ induces an action of the abstract group $\ti\mT$ on $\ti\mV$.

\end{enumerate}

\end{lemma}

\begin{proof}

At the finite level this is the standard currying isomorphism $\operatorname{Maps}(X\times Y,V)\simeq\operatorname{Maps}(X,\operatorname{Maps}(Y,V))$, which is compatible with the one-point set and with the usual coherence diagrams. Passing from finite sets to locally profinite objects replaces maps by locally constant  supported functions. The same currying map identifies locally constant compactly supported functions on $\mathbf X\times\mathbf Y$ with locally constant  supported families, parametrized by $\mathbf X$, of locally constant  supported functions on $\mathbf Y$. Thus $a_1(\mathbf X\times\mathbf Y,V)\simeq a_1(\mathbf X,a_1(\mathbf Y,V))$, and the unit object acts as the identity.

\sms

The functors $a_2$, $a_3$, and $a$ are obtained from $a_1$ by the countable ind- and pro-limits appearing in the definitions. Products and all structure maps are represented levelwise, so the preceding associativity and unit isomorphisms pass to these limits. Since the coherence diagrams already commute at the finite level, they commute after passing to the ind/pro limits. This proves that $a$ is an action of the monoidal category $\mSet$ on $\mVect$ in the stated convention.

\sms

For the second assertion, let a group object $\mT$ act on $\mV$, so that we have an action morphism $\mV\to a(\mT,\mV)$. Applying the functors $\ti{\ }$ identifies this morphism with a map $\ti\mV\to\operatorname{Maps}(\ti\mT,\ti\mV)$, or equivalently with a map $\ti\mT\times\ti\mV\to\ti\mV$. The associativity and unit identities follow by applying $\ti{\ }$ to the corresponding diagrams for the action of $\mT$ on $\mV$. This is the required compatibility.

\end{proof}

\begin{remark}

This action induces the pseudo-action described in Section $1.9$ of \cite{GK}.

\end{remark}

\subsubsection{Quasi-pro-unipotent group objects, coinvariants, and inflation}

\begin{definition}\label{quasi-pro-unipotent}
\begin{enumerate}
\item A group object $H\in\bSet$ is quasi-unipotent if it admits a presentation
$$
H\simeq\inl_i H_i,
$$
where every $H_i$ is a group object of $\Pro_{\aleph_0}(Set_0)$ and all transition morphisms are homomorphisms.
\item A group object $H\in\Pro\bSet$ is quasi-pro-unipotent if it admits a presentation
$$
H\simeq\prl_j H_j,
$$
where every $H_j$ is quasi-unipotent and the transition homomorphisms are  surjective.
\end{enumerate}
\end{definition}

Let $H$ be a group object for which the coinvariants functor is defined. We denote it by
$$
\operatorname{Coinv}_H:\Rep(H,\mVect)\longrightarrow\mVect,
$$
and write $\Pi_H:=\operatorname{Coinv}_H(\Pi)$. Thus $\operatorname{Coinv}_H$ is left adjoint to the trivial-representation functor $\operatorname{triv}:\mVect\to\Rep(H,\mVect)$.

\begin{theorem}\label{coinvariants-exact}
If $H$ is quasi-pro-unipotent, then the functor $\operatorname{Coinv}_H$ is exact.
\end{theorem}

\begin{proof}
This is \cite[Lemma~2.7]{GK}. Its hypotheses are precisely that $H$ be an inverse limit, with surjective transition homomorphisms, of group objects which are direct limits of compact group objects.
\end{proof}

\begin{proposition}\label{inflation}
If $H$ is quasi-pro-unipotent, then $\operatorname{Coinv}_H$ admits a left adjoint
$$
\operatorname{Inf}_H:\mVect\longrightarrow\Rep(H,\mVect),
$$
called inflation. Thus, functorially in $W\in\mVect$ and $\Pi\in\Rep(H,\mVect)$,
$$
\Hom_{\Rep(H,\mVect)}\bigl(\operatorname{Inf}_H(W),\Pi\bigr)
\simeq
\Hom_{\mVect}\bigl(W,\Pi_H\bigr).
$$
\end{proposition}

\begin{proof}
This is \cite[Proposition~2.7]{GK1}. Notice that this proposition, rather than the exactness statement in Theorem \ref{coinvariants-exact}, is the source of representability. In particular,
$$
\Hom_{\Rep(H,\mVect)}\bigl(\operatorname{Inf}_H(\CC),\Pi\bigr)
\simeq
\Hom_{\mVect}(\CC,\Pi_H).
$$
Consequently the functor $\Pi\mapsto\Hom_{\mVect}(\CC,\Pi_H)$ is represented by $\operatorname{Inf}_H(\CC)$. Whenever $\Pi_H$ is an ordinary vector space, the right-hand side is canonically identified with $\Pi_H$.
\end{proof}

Assume now that $F$ is a non-Archimedean local field with ring of integers $O_F$ and uniformizer $\varpi$, and put $O=F[[t]]$.

\begin{lemma}\label{tnO-quasi-pro-unipotent}
For every $n\geq0$, the additive group object $t^{-n}O$ is quasi-pro-unipotent. Consequently the functor of $t^{-n}O$-coinvariants is exact, and it admits the left adjoint $\operatorname{Inf}_{t^{-n}O}$.
\end{lemma}

\begin{proof}
For $r\geq1$, set
$$
H_r:=t^{-n}O/t^rO.
$$
Then
$$
t^{-n}O\simeq\prl_{r\geq1}H_r,
$$
and every transition homomorphism $H_{r+1}\to H_r$ is surjective. Moreover,
$$
H_r\simeq F^{n+r}
$$
as an additive locally profinite group, and
$$
F^{n+r}=\inl_{s\geq0}\varpi^{-s}O_F^{n+r}.
$$
Each group $\varpi^{-s}O_F^{n+r}$ is compact and profinite, hence is a group object of $\Pro_{\aleph_0}(Set_0)$. Therefore every $H_r$ is quasi-unipotent and $t^{-n}O$ is quasi-pro-unipotent. Exactness follows from Theorem \ref{coinvariants-exact}, and the existence of inflation follows from Proposition \ref{inflation}.
\end{proof}

\subsubsection{Group objects $\mH (K)$}

As is explained in section $2.11$ of \cite{GK} with any  algebraic
$F$-group $ H$ one can associate a group object
 $\mH\in \mSet$ such that $\wt {\mH} $ is equal (as an abstract group) to $H(K)$.

The following statement is Lemma $2.13$ in \cite{GK}.

\begin{claim}

The functor $\text{rest}_{\mH(K)}:\mcR(\mH)\to\Rep(\ti\mH)$ is fully faithful and defines an equivalence between the category $\mcR(\mH)$ and the category $\mcR (\ti \mH)$ of smooth representations of the group $\ti \mH$.
\end{claim}

Thus we shall freely pass between representations of $\mH$ (which is a group object in the category $\mSet$) and smooth representations of the abstract group $\ti\mH$ (let us again stress that in both cases we are talking about representations on objects of $\mVect$).

\begin{definition} Let $\mT $ be a group object in $ \textbf{Set} , (\Pi, \mV)$ a smooth representation of the group $\ti \mT$ and $\kk : \ti \mT \to \mC ^*$ a smooth character. We define $\mcF_\kk:\mVect\to Set$ to be the functor given by $\mW \to \{ \beta \in \Hom _{\mVect}(\mV ,\mW)| \beta \circ \ti t = \kk (\ti t)\beta , \ti t \in \ti \mT  \}$.
\end{definition}
The following representability statement follows from the inflation adjunction; see Proposition \ref{inflation}.
\begin{claim}\label{coin} The functor $ \mcF _\kk $ is representable by an object which we denote by $\mV _{\mT ,\kk}$. We write $\mV _\mT$ instead of $\mV _{\mT ,1}$.
\end{claim}
The following statement is implied by Proposition $2.5$ in \cite{GK}.
\begin{claim}\label{2.5}
Any representation of a strict group object $\mH\in \Pro\textbf{Set}$ on the category $\mVect$
is a projective limit of representations of $\mH$ on the category $\Vect _ {\aleph _0} $.
\end{claim}

If $\mH\in \mSet$ is a group object and $\mL\subset\mH$ is a group subobject, we have the restriction functor $r^\mH _\mL$. The following statement is Proposition $3.5$ of \cite{GK}.
\begin{claim}\label{indH}If $\mH /\mL$ is ind-compact, the functor $i_\mH^\mL$ of induction is the right adjoint to the restriction functor $r^\mH_\mL$.
\end{claim}

\subsection{Examples of group objects in $  \mSet$}
\subsubsection{The group object
$\mP$}

\begin{definition}\label{P_0}

 \begin{enumerate}

\item
We denote by $\ti v:P(K) \to \mZ$ the
morphism such that $\ti v \begin{pmatrix}a&b\\0&1 \end{pmatrix}= v(a)$
and write $P(K)_0:=\ker(\ti v)$.
So $P_0:= O^* \ltimes K \subset P(K)$.
\item $  P ^n(K)=\{(a,b)\in K^* \ltimes K | \min(v(a),-v(a),v(b))\geq -n\}$.

\item
$\GG _m \subset P(O) $  is the subgroup  of pairs  $(a,b)\in O^\ast \times O$ such
that $v(a-1),v(b)\geq m$. \end{enumerate} \end{definition}
The subsets $P ^n(K)\subset  P$ are two-sided $P(O)$-invariant and we write  $ P ^n _m := P^n (K)/ \GG_m $.
\begin{claim} $P ^n _m$ are
$F$-points of smooth algebraic $F$-varieties.
\end{claim}
\begin{definition}

 \begin{enumerate}
\item Since $P^n_m$ are
$F$-points of smooth algebraic $F$-varieties they define objects $\mP ^n _m \in  \bSet$.

\item  $\mP ^n := \prl _m \mP ^n_m \in \Pro \textbf{Set}$ where $ \mP ^n _{m+1} \to \mP ^n _m $ are the natural  smooth morphisms.
\item    $\mP := \inl _n \mP ^n\in \mSet$.
\end{enumerate} \end{definition}
\begin{claim}
$ \mP \in \mSet $ has a natural structure of group
object and   $\wt {\mP}= P(K)$.
\end{claim}

\begin{remark} We have $\mP = \mK ^\ast \ltimes \mK$ where 
$\mK ^\ast ,\mK $ are group objects such that  $\wt {\mK^\ast}= K\ast$ and  $\wt {\mK}= K$.
\end{remark}

\subsubsection{The group object
$ \mG$}

Let $ G(K):=  {\PGL}_2(K)$.  \begin{definition}\label{G}
\begin{enumerate}
\item $ \hat  X ^n\subset \text{GL}_2(K),n\geq 1$ is  the set of $2\times 2$-matrices $\{a_{ij}\}$
such that
 $2 \min \limits_{1\leq i,j \leq 2}v(a_{ij})- v(\det (A))\geq -n$.
\item $  X ^n\subset G(K),n\geq 1$ is the image of $  \hat  X^n$ under the projection $ \text{GL}_2(K) \to  G(K)$.

\item $ X ^n  _m:=  X^n /  G(O)(m+1) $ (see Definition \ref{G-prime-filtration}).

\sms

  $X ^n _m$ are
$F$-points of algebraic $F$-varieties and we  denote by  $ \mX ^n _m $ the corresponding objects in $  \textbf{S} et$.
\item
$\mX ^n := \prl _m X^n_m \in \Pro \textbf{Set}$ where $ X ^n _{m+1} \to X ^n _m $ are the natural  smooth morphisms.
\item
 $\mG:= \inl _n \mX ^n\in \mSet$.
\end{enumerate}
\end{definition}
\begin{claim} $ \mG \in \mSet $ has a natural structure of group object and $\wt {\mG}=G(K)$.
\end{claim}

\subsection{Smooth irreducible representations of the group $P(O)$}

\begin{lemma}\label{inf}
Every smooth irreducible representation of $P(O)$ on $\mVect$ is the inflation of a representation of $P(O)_m$ for some $m\geq0$.
\end{lemma}
\begin{proof}
By \cite[Proposition~2.5]{GK}, every smooth representation of the strict group object $P(O)$ on $\mVect$ admits a strict presentation
$$
\mV\simeq\varprojlim_r V_r
$$
by smooth representations on ordinary vector spaces, with surjective transition morphisms. Let $\mV$ be nonzero and irreducible. Choose $r$ such that the canonical projection $p_r:\mV\to V_r$ is nonzero. Its kernel is a $P(O)$-stable subobject of $\mV$ and hence is zero. Since $p_r$ is also an epimorphism, it is an isomorphism. Thus $\mV$ is an ordinary smooth irreducible representation of $P(O)$.
Choose $0\neq v\in\mV$. Smoothness in the sense of \cite[Subsection~2.4]{GK} implies that the $P(O)$-subrepresentation generated by $v$ factors through $P(O)/\GG_m$ for some $m\geq0$. This subrepresentation is nonzero and therefore, by irreducibility, is all of $\mV$. Hence $\GG_m$ acts trivially on $\mV$, and $\mV$ is the inflation of an irreducible representation of $P(O)_m=P(O)/\GG_m$.
\end{proof}
\begin{corollary}
The definition of smoothness for irreducible representations of $P(O)$ used here agrees with the definition in \cite[Subsection~2.4]{GK}.
\end{corollary}
\subsection{A representation $(\rho_m,\mcM),m\in \mZ$ of $P_0=O^*\ltimes K$}
For $n\geq0$, put
$$
Z_n=O^*/O^*(n+1)
$$

and let $\mcM_n=\operatorname{Meas}_c(Z_n)$ be the space of locally constant compactly supported complex measures on $Z_n$. Pushforward along the natural projection $Z_{n+1}\to Z_n$ gives a surjective map
$$
q_n:\mcM_{n+1}\longrightarrow\mcM_n.
$$
\begin{definition}\label{M-definition}
Put
$$
\mcM=\varprojlim_{n\geq0}\mcM_n.
$$
There is a representation $\rho_0:P_0\to\Aut(\mcM)$ characterized as follows. For $a\in O^*$, the element $\widehat a$ acts by pushforward under $z\mapsto az$. For $b\in K$, the element $\widetilde b$ acts by multiplication by the function
$$
z\longmapsto\psi(bz/t).
$$
More precisely, if $b\in t^{-n}O$, this function is well defined on $Z_n$, and these finite-level actions are compatible with the transition maps. Since every $b\in K$ belongs to $t^{-n}O$ for some $n$, they define an action of $K$ on the inverse system $\mcM$.
\end{definition}
\begin{claim}\label{M-smooth-coinvariants}
\begin{enumerate}
\item The representation $(\rho_0,\mcM)$ is smooth.
\item For every $n\geq0$, the canonical projection $\mcM\to\mcM_n$ induces an isomorphism
$$
\mcM_{O^*(n+1)}\stackrel{\sim}{\longrightarrow}\mcM_n.
$$
\end{enumerate}
\end{claim}
\begin{proof}
The first assertion follows from the finite-level description in Definition \ref{M-definition}. For the second, $O^*(n+1)$ acts trivially on $Z_n$. At every higher level, taking $O^*(n+1)$-coinvariants identifies measures in the fibers of the projection to $Z_n$; the resulting map is pushforward. Passing to the projective limit gives the asserted isomorphism.
\end{proof}
For $n\geq0$, put
$$
U^n=\left\{\widetilde b:b\in t^{-n}O\right\},\qquad \overline U^n=U^n/U^{-1},
$$
where $U^{-1}=\{\widetilde b:b\in tO\}$. The action of $O^*\ltimes U^n$ on $\mcM_n$ factors through the algebraic $F$-group
$$
P_0^n=(O^*/O^*(n+1))\ltimes\overline U^n.
$$
Denote this finite-level representation by $\rho_0^n$.
\begin{lemma}\label{M-irred}
The representation $(\rho_0,\mcM)$ is irreducible.
\end{lemma}
\begin{proof}
The normal commutative subgroup $\overline U^n$ acts on $\mcM_n$ with character support
$$
\left\{b\longmapsto\psi(bx/t):x\in X_n\right\}.
$$
The group $O^*/O^*(n+1)$ acts simply transitively on this set. The little-group theorem therefore shows that $\rho_0^n$ is irreducible.
Let $\mN\subset\mcM$ be a nonzero $P_0$-stable subobject, and let $N_n$ be its image in $\mcM_n$. For some $n$, one has $N_n\neq0$, and irreducibility gives $N_n=\mcM_n$. For $r\geq n$, surjectivity of $q_{r-1}\circ\cdots\circ q_n$ implies $N_r\neq0$, so $N_r=\mcM_r$ by irreducibility. It follows also that $N_r=\mcM_r$ for $r<n$. Hence $\mN\to\mcM_n$ is surjective for every $n$. Since $\mcM$ is a strict pro-object, the strictness criterion for morphisms of pro-vector spaces implies $\mN=\mcM$.
\end{proof}
\begin{definition}
For $m\in\ZZ$, define a representation $(\rho_m,\mcM)$ of $P_0$ by
$$
\rho_m(p):=\rho_0(\widehat t^{-m}p\widehat t^m),\qquad p\in P_0.
$$
Put
$$
\Upsilon:=\prod_{m\in\ZZ}\mcM.
$$
The group $P_0$ acts on the $m$-th factor by $\rho_m$, and $\widehat t$ acts by the shift from the $m$-th factor to the $(m+1)$-st factor. These actions define a representation of $P(K)=\langle\widehat t\rangle\ltimes P_0$ on $\Upsilon$.
\end{definition}

\subsection{Identification with induction from $P(O)$}
Put
$$
P'(O)=O^*(1)\ltimes O
$$
and let $\widehat\psi$ be the character of $P'(O)$ which is trivial on $O^*(1)$ and whose restriction to $O$ is
$$
\psi_0(c)=\psi(c/t).
$$
As follows from Lemma \ref{sec}, we have
$$
\bs(0)\simeq\ind_{P'(O)}^{P(O)}\widehat\psi.
$$
Let
$$
(\tau_m,\Upsilon_m)=\ind_{P(O)}^{P_0}(\bs(m),\mcS(m)).
$$

\begin{proposition}\label{Upsilonm-M}
There is a canonical isomorphism of $P_0$-representations
$$
(\tau_m,\Upsilon_m)\simeq(\rho_m,\mcM),\qquad m\geq0.
$$
\end{proposition}
\begin{proof}
We first treat $m=0$.

\sms

Induction in stages gives
$$
\Upsilon_0\simeq\ind_{P'(O)}^{P_0}\widehat\psi.
$$
For $n\geq0$, put $P_0^{\leq n}=O^*\ltimes t^{-n}O$ and let
$$
\Upsilon_0^n=\ind_{P'(O)}^{P_0^{\leq n}}\widehat\psi.
$$
Then
$$
\Upsilon_0\simeq\varprojlim_{n\geq0}\Upsilon_0^n.
$$
As an $F$-variety, the quotient $P_0^{\leq n}/P'(O)$ admits coordinates
$$
F^*\times(t^{-n}O/O).
$$
The character $\widehat\psi$ supplies the usual twisted translation action on the second factor. The residue pairing
$$
(t^{-n}O/O)\times(O/t^nO)\longrightarrow F,\qquad
(z,y)\longmapsto\operatorname{Res}_{t=0}(zy\,dt)
$$
is perfect. Fourier transformation in the second variable therefore gives an isomorphism
$$
\mathcal F_n:\Upsilon_0^n\stackrel{\sim}{\longrightarrow}
\operatorname{Meas}_c\bigl(F^*\times(O/t^nO)\bigr).
$$
The map
$$
F^*\times(O/t^nO)\longrightarrow O^*/O^*(n+1),\qquad
(a,y)\longmapsto a(1+ty),
$$
is an isomorphism of $F$-varieties. Under this isomorphism, the action of $\widehat a\in O^*$ is pushforward under $x\mapsto ax$, while the action of $\widetilde b$, $b\in t^{-n}O$, is multiplication by
$$
x\longmapsto\psi(bx/t).
$$
Thus $\mathcal F_n$ identifies $\Upsilon_0^n$ with $\mcM_n$ as a representation of $P_0^{\leq n}$.
Choose the self-dual additive Haar measures in the finite Fourier transforms. Then restriction from $t^{-n-1}O/O$ to $t^{-n}O/O$ is transformed into pushforward along
$$
O^*/O^*(n+2)\longrightarrow O^*/O^*(n+1).
$$
Consequently the isomorphisms $\mathcal F_n$ are compatible with the transition maps. Passing to projective limits gives
$$
\Upsilon_0\simeq\varprojlim_n\mcM_n=\mcM
$$
as $P_0$-representations.

For general $m\geq0$, Lemma \ref{sec} and induction in stages give
$$
\Upsilon_m\simeq\ind_{\widetilde R_m}^{P_0}(\bl_m),
\qquad \widetilde R_m=O^*(1)\ltimes t^mO.
$$
Conjugation by $\widehat t^{-m}$ identifies $\widetilde R_m$ with $P'(O)$. If $c\in O$ and $b=t^mc\in t^mO$, then
$$
\bl_m(1,b)=\psi(t^{-m-1}b)=\psi(c/t)=\widehat\psi(1,c).
$$
Thus this conjugation identifies the inducing pair $(\widetilde R_m,\bl_m)$ with $(P'(O),\widehat\psi)$. Consequently it identifies the resulting $P_0$-action with the twist
$$
p\longmapsto\rho_0(\widehat t^{-m}p\widehat t^m)=\rho_m(p).
$$
Hence $\Upsilon_m\simeq\mcM$ with the action $\rho_m$, as asserted.
\end{proof}

\section{Smooth representations of $P(K)$}\label{P(K)}

\subsection{Smooth representations of the group $K$}
\begin{definition} $\mK$ and $\mK ^*$ are group objects of $\mSet$ which are considered as subgroups of $\mP$.
\end{definition}
So $K$ and $K^*$ are considered as subgroups of $P(K)$.

 \begin{definition}\label{cusp} Let $\mW$ be a smooth representation of $K$.
\begin{enumerate}
\item
$\eta _n (\mW):\mW \to \mW _ {t^nO}$ is the canonical projection (see Claim \ref{coin}).
\item $\mcF _\mW= \{\mF _n\}$ is a decreasing filtration on $\mW$ where
$ \mF _n := \ker (\eta _n(\mW))$.
\item
$\kk ^n_m(\mW):   \mW _ {t^nO} \to  \mW _ {t^mO}, m\leq n\in \mZ $ are the canonical projections.
\item $ \mW _n^m:= \ker ( \kk ^n_m(\mW)) $.
\item $\mu _n(\mW) :=  \mW ^{n-1}_n $.

\item $\mu (\mW):=  \mu _0 (\mW)$.
\end{enumerate} \end{definition}

\begin{lemma}\label{m} \begin{enumerate}

\item The functors $\eta_n$ are exact.

\item We have a natural smooth representation of the group $K$ on $\mu (\mW)$ trivial on $O$.
\end{enumerate}
\end{lemma}
\begin{proof} Part (1) follows from Lemma 2.7 in \cite{GK}, and part (2) follows from part (1).
\end{proof}

 \subsection{The exactness of the functor $\mu$}

\begin{proposition}\label{mmu} The functor $\mu$ is exact.
\end{proposition}
\begin{proof} The proof of Proposition \ref{mmu} is based on the following result.

\sms

Let $\mcC$ be an abelian category, let $F,G:\mcC\to\mcC$ be exact functors, and let $a:F\to G$ be a morphism such that the maps $a(X):F(X)\to G(X)$ are surjective for all $X\in Ob(\mcC)$. We define $R : = \ker (a)$.
\begin{lemma}\label{exact}The functor $R$ is exact.
\end{lemma}
\begin{proof}
Let
$0\to A'\xrightarrow{\,i\,} A\xrightarrow{\,p\,} A''\to0$
be a short exact sequence in $\mcC$. Since the functors $F$ and $G$ are exact we get two short exact rows and a commutative diagram

$\begin{array}{ccccccccc}
0&\to&F(A')&\xrightarrow{F(i)}&F(A)&\xrightarrow{F(p)}&F(A'')&\to&0\\
& &\downarrow^{a_{A'}}&&\downarrow^{a_A}&&\downarrow^{a_{A''}}&&\\
0&\to&G(A')&\xrightarrow{G(i)}&G(A)&\xrightarrow{G(p)}&G(A'')&\to&0.
\end{array}$

The snake lemma provides  the exact sequence
$$0\to\ker(a_{A'})\to\ker(a_A)\to\ker(a_{A''})\to\operatorname{coker}(a_{A'})\to\cdots$$
We have $\operatorname{coker}(a_X)=0$ since, by hypothesis, each $a_X:F(X)\to G(X)$ is surjective.  Therefore the snake lemma yields the exactness of the sequence

$0\to\ker(a_{A'})\to\ker(a_A)\to\ker(a_{A''})\to0$. Thus $R$ is exact.
\end{proof}
Proposition \ref{mmu} now follows from Lemma \ref{m}.

\end{proof}

\subsection{Representations of  $P(K)$} 
We now introduce some basic constructions related to representations of the group $P(K)$ that we are going to use in the future.

\begin{definition}\label{grad}
\begin{enumerate}
\item A smooth representation of $P(K)$ is complete if its filtration $\mcF_{\mW}$ is admissible in the sense of
Definition \ref{admissible}.
\item 
For  a smooth representation $(\Pi ,\mW)$ of $P(K)$ we write  $\bar \mW := \prod _{n\in \mZ}
\mu _n(\mW)$. Since the group $P(K)$ preserves (up to a shift) the filtration $\mcF$ on $\mW$ we obtain a representation $\bar\Pi$ of the group $P(K)$ on $\bar \mW$. 
\end{enumerate}
\end{definition}
\begin{claim}
Let $\Pi$ be a smooth  representation of $P(K)$. Then the representation $\bar\Pi$ is also smooth.
\end{claim}

\begin{proposition}\label{isbar}

Let $\mW$ be a smooth complete representation of $P(K)$ such that $\bar \mW$ is isomorphic to $\Upsilon$ and
 $\mW ^{N(K)}=\{0\}$. Then $\mW$ is irreducible.

\end{proposition}
\begin{proof}Let $\mU\subset\mW$ be a nonzero $P(K)$-invariant subobject. Apply Lemma \ref{bar} to the increasing filtration $\mW_i:=\mF_{-i}$. It is sufficient to show that $\mU\cap\ker(\eta_n)\neq0$ for some $n\in\ZZ$. We show that the assumption
$$
\mU\cap\ker(\eta_n)=0\qquad(n\in\ZZ)
$$
leads to a contradiction.

\sms

It is clear that $\operatorname{Im}(\operatorname{Id}-\pi(\widetilde b))\subset\ker(\eta_n)$ for $n<v(b)$. Therefore the assumption above implies that $N(K)$ acts trivially on $\mU$. Hence $\mU\subset\mW^{N(K)}=0$, a contradiction.
\end{proof}
%--------------------------------------------------------------------------------------------------------

\section{The affine Grassmannian}\label{Gr}
In what follows we would like to study induction of special representations of $G(O)$ to $G(K)$. In order to do this we need to discuss some preparatory material about the geometry of the affine Grassmannian $\Gr_G=G(K)/G(O)$. We refer the reader to Sections 2 and 3 of \cite{MV} for the proof of all the claims about $\Gr_G$ which are going to appear in this Section (and which are all easy in any case).
Similarly, in order to analyze the induction of special representations of $G'(O)$ to $G(K)$ we need to discuss the twisted affine Grassmannian $\Gr_G'=G(K)/G'(O)$; this is done at the end of this section. We should warn the reader that in this section we abuse notation and usually do not distinguish between an ind-scheme over $F$ and its set of $F$-points, considered as a topological space with respect to the topology induced by that on $F$; we hope that this will not cause confusion.
%--------------------------------------------------------------------------------------------------------------------------
\subsection{Topology on the affine Grassmannian}
For any algebraic group $H$ over $F$ we set $\Gr_H=H(K)/H(O)$. It is well known that this is the set of $F$-points of an ind-finite-dimensional ind-scheme over $F$; in particular $\Gr_H$ (as a set) can be written as an inductive limit of finite-dimensional $F$-varieties with respect to closed embeddings. This endows the set $\Gr_H$ with a natural topology.

\begin{remark}The topological space $\Gr_H$ is usually not locally compact.
\end{remark}
\subsection{The basic geometry of $\Gr_G$}
The following statement follows from the Iwasawa decomposition proved in \cite{T}.
\begin{claim}\label{trans} The group $P(K)$ acts transitively on $\Gr_G$. Equivalently, the natural map $j:\Gr_P\to \Gr_G$ is a bijection.
\end{claim}
\begin{remark}\label{not-isom}
Although the map $j$ is a bijection, it is not an isomorphism either on the level of ind-schemes or on the level of topological spaces (the latter just means that the inverse map is not continuous). In fact, it is observed in \cite{BD} (cf. the lemma at the end of page 173) that, for any affine algebraic group $G$ and any closed subgroup $H$, the natural map $\Gr_H\to\Gr_G$ is a locally closed embedding of ind-schemes if the quotient $G/H$ is a quasi-affine variety (in particular, in this case $\Gr_H$ is isomorphic to its image in $\Gr_G$). However, in our case the quotient $G/P$ is the projective line $\mathbb P^1$, which is not quasi-affine, so the above lemma is not applicable.
\end{remark}
For an integer $l\geq 0$, let $\Gr_G^l$ denote the $G(O)$-orbit of the image in $\Gr_G$ of the diagonal matrix
$\hat t\, ^l=\begin{pmatrix} t^l & 0\\ 0 & 1\end{pmatrix}$.  Then $\Gr_G^l$ is an algebraic $F$-variety of dimension $l$. We denote by $\oGr_G^l$ (the set of $F$-points of) its Zariski closure. This is a projective $F$-variety of dimension $l$. Moreover, we have
\begin{equation}
\Gr_G=\bigsqcup\limits_{l=0}^{\infty} \Gr_G^l;\ \ \  \oGr_G^l=\bigsqcup\limits_{0\leq i\leq l, \ \text{$i-l$ is even}} \Gr_G^i;
\end{equation}
In particular, we have closed embeddings $\oGr_G^l\hookrightarrow \oGr_G^{l+2}$, and $\Gr_G$ is equal to the inductive limit of all the $\oGr_G^l$ (note that as an ind-scheme it has two connected components corresponding to the parity of $l$).

\subsection{Semi-infinite orbits}
The Iwasawa decomposition implies that
\begin{equation}\label{semi-inf}
\Gr_G=\bigsqcup\limits_{n\in \mathbb Z} N(K)\cdot \hat t\, ^n,
\end{equation}
where as before we do not distinguish notationally the matrix $\hat t\, ^i$ and its image in $\Gr_G$. We set $S^n=N(K)\cdot \hat t\, ^n$ and call these sets "semi-infinite orbits" (associated with the choice of the maximal unipotent subgroup $N$). Each $S^n$ is actually invariant under $P_0=T(O)N(K)$ and it is a locally closed sub-ind-scheme of $\Gr_G$.
Note that the stabilizer of $\hat t^i$ in $N(K)=K$ is equal to $t^i O$, thus $S^i$ is isomorphic to $K/t^i O$. It is easy to see that this is an isomorphism of ind-schemes over $F$ and thus also of topological spaces (this is in fact a special case of the Lemma at the end of page 173 of \cite{BD}). The closure of $S^i$ (in either Zariski topology or the topology induced by the topology on $F$) is equal to the union of all the $S^j$ with $j\leq i$.
Note also that $\Gr_P$ is isomorphic to the disjoint union of the $S^i$.
\footnote{Warning: this statement is true for topological spaces with respect to the ``classical'' topology (i.e., the topology induced by that on $F$), or for reduced ind-schemes. However, it is not true for ind-schemes themselves: in fact, $\Gr_P$ is highly nonreduced. This subtlety will not be important below.}
The intersection $S^i\cap \Gr_G^l$ is non-empty if and only if $|i|\leq l$  and $l-i$ is even, and in this case this intersection is isomorphic to $\mathbb A^{\frac{l+i}{2}}$ (as an algebraic variety over $F$).

\subsection{The ind-schemes $\Gr_{G,m}$} Let us now consider the quotient $\Gr_{G,m}=G(K)/G(O)(m+1)$. This is the set of $F$-points of an ind-scheme over $F$ which has the structure of a principal $G(O)_m$-bundle over $\Gr_G$. Explicitly, it is equal to the union of finite-dimensional schemes $\overline\Gr_{G,m}^l$ where $\overline\Gr_{G,m}^l$ is the preimage of
$\overline\Gr_G^l$ under the natural projection $\Gr_{G,m}\to \Gr_G$.

The ind-scheme $\Gr_{G,m}$ makes sense for any algebraic group, so in particular, we can consider $\Gr_{P,m}$. We have an obvious morphism $\Gr_{P,m}\to \Gr_{G,m}$.
The natural map
$$
\Gr_{P,m}\to \Gr_P\underset{\Gr_G}\times \Gr_{G,m}
$$
is a closed embedding of ind-schemes (since the RHS is a $G(O)_m$-torsor over $\Gr_P$ and the LHS is a $P(O)_m$-sub-torsor (and $P(O)_m$ is closed in $G(O)_m$)). This implies the following. Since the group ind-scheme $P(K)$ is equal to the union of $\mathbb Z$ connected components (corresponding to the homomorphism $P(K)\to \mathbb Z$ given by the composition of $v$ with the natural map $P(K)\to K^*$), the same is true for each of the ind-schemes $\Gr_{P,m}$. Let $\Gr_{P,m}^i$ denote the corresponding $i$-th connected component (here we use the Zariski topology, not the topology coming from the topology on $F$). Note that $\Gr_{P,0}^i=S^i$ and we shall sometimes also denote it by $\Gr_P^i$. On the other hand, on the level of reduced ind-schemes $\Gr_P\underset{\Gr_G}\times \Gr_{G,m}$ is isomorphic to the disjoint union of $S^i_m:=p_m^{-1}(S^i)$ where $p_m:\Gr_{G,m}\to \Gr_G$ is the natural projection. Thus we see that
the reduced ind-scheme of $\Gr_{P,m}^i$ is closed inside $S^i_m$.
\begin{corollary}
The topological space $\Gr_{P,m}^i$ is closed in $S^i_m$.
\end{corollary}

\subsection{The twisted affine Grassmannian}We now define
$$
\Gr'_G=G(K)/G'(O).
$$
 Again, this is the set of $F$-points of an ind-finite-dimensional ind-proper ind-scheme over $F$. It is still true that $G(K)=P(K)\cdot G'(O)$ (cf. Lemma \ref{BT}). Hence as before we can set $(S')^i$ to be the (isomorphic) image of $\Gr_P^i$ in $\Gr'_G$.

For $m\in \mathbb Z_{\geq 0}$ we can again set $\Gr'_{G,m}=G(K)/G'(O)(m+1)$; this is a $G'(O)_m$-torsor over $\Gr'_G$ and similarly to the above we define $(S'_m)^i$ to be the preimage of $(S')^i$. It contains $\Gr_{P,m}^i$ as a closed subset.

\section{Special representations of $G(K)$ -- some initial results}\label{G(K)-in}

\subsection{Restriction from $G(K)$ to $P(K)$} We would like to study representations of $G(K)$ whose restriction to $P(K)$ is irreducible. Recall that if we replace $K$ by $F$, then this property is equivalent to irreducibility and cuspidality. Moreover, the group $P(F)$ has a unique (up to isomorphism) irreducible cuspidal representation, naturally realized in the space $\mcS(F^*)$. When $F$ is replaced by $K$, the uniqueness statement is no longer true. The most straightforward analogue of the unique irreducible cuspidal representation of $P(F)$ is the representation $\Upsilon$ of $P(K)$ defined above. It is therefore natural to ask whether there are irreducible representations of $G(K)$ whose restriction to $P(K)$ is isomorphic to $\Upsilon$. In this section we show that the answer is ``no.'' However, we shall see later that, if $\Pi(\pi)$ is induced from a special representation of either $G(O)$ or $G'(O)$ and $\Theta(\pi):=\Pi(\pi)|_{P(K)}$, then $\Theta(\pi)$ is complete, has no $N(K)$-invariants, and satisfies
$$
\bar\Theta(\pi)\simeq\Upsilon.
$$
Proposition \ref{isbar} will then imply that $\Theta(\pi)$ is irreducible.

As before we identify $K$ with $N(K)\subset P(K)$ and therefore identify $O$
 with a subgroup of $G(K)$. We define $G^+(K)\subset G(K)$ to be the image of $\SL_2(K)$ under the natural map $\SL_2(K)\to\PGL_2(K)$.

\begin{lemma}\label{Gplus-invariants}
Let $(\Pi,\mW)$ be a smooth representation of $G(K)$. Then the natural inclusion
$$
\mW^{G^+(K)}\hookrightarrow\mW^O
$$
is an isomorphism.
\end{lemma}
\begin{proof}
The subobject $\mW^O$ is invariant under $N(K)$ because $N(K)$ is abelian. Let $b\in K$. By Proposition \ref{infty}, applied to the restriction to $G^+(O)$, every vector of $\mW^O$ is fixed by $G^+(O)$. Applying this to $\widetilde b\,v\in\mW^O$ shows that $v$ is fixed by
$$
\widetilde b^{-1}G^+(O)\widetilde b.
$$
As $b$ varies, these conjugates generate the image of $\SL_2(K)$ in $\PGL_2(K)$. By definition this image is $G^+(K)$. Thus $\mW^O\subset\mW^{G^+(K)}$. The reverse inclusion is immediate because $O\subset G^+(K)$.
\end{proof}
\begin{corollary}\label{no-rest}There is no representation of $G(K)$ whose restriction to $P(K)$ is isomorphic to $\Upsilon$.
\end{corollary}
\begin{proof}
Suppose that a representation $(\Pi,\mW)$ of $G(K)$ satisfies $\mW|_{P(K)}\simeq\Upsilon$. The space $\Upsilon^O$ is nonzero: for example, in the realization of Definition \ref{M-definition}, the subgroup $O$ acts trivially on every factor $(\rho_m,\mcM)$ with $m<0$. Hence $\mW^O\neq0$. By Lemma \ref{Gplus-invariants},
$$
\mW^O=\mW^{G^+(K)}.
$$
The group $G^+(K)$ is normal in $G(K)$, so $\mW^{G^+(K)}$ is a nonzero $G(K)$-invariant, and hence $P(K)$-invariant, subobject of $\mW$. Since $\Upsilon$ is irreducible, it follows that $\mW^{G^+(K)}=\mW$. Thus $N(K)\subset G^+(K)$ acts trivially on $\mW$, contradicting the nontrivial action of $N(K)$ on $\Upsilon$.
\end{proof}

\section{Representations of $G(K)$ induced from special representations of $G(O)$}\label{GKinduced}
In this section we study representations of $G(K)$ induced from special representations of $G(O)$. We freely use the notation of Section \ref{Gr}.

\subsection{Induction from $H(O)$ to $H(K)$ for general algebraic group $H$}\label{ind-gen}
We would like to discuss the functor of compact induction $\ind_{H(O)}^{H(K)}$ for a general linear algebraic group $H$. This functor is well-defined but for general $H$ it is not adjoint to the restriction functor on any side. However, it is actually right adjoint to the restriction functor if $H$ is reductive. The definition of this functor, as well as the proof of the above statement can be found in Section 3 of \cite{GK}. For the reader's convenience we reproduce the definition here.

Namely, let $(\pi,V)$ be a smooth representation of $H(O)$. Since such a $\pi$ is always equal to the projective limit of representations for which $V$ is a vector space (as opposed to a pro-vector space) and since the functor $\ind$ will commute with such limits, it is enough to assume that $V$ is a vector space. Such a $\pi$ is the inductive limit of representations trivial on $H(O)(m+1)$ for some $m$, and since $\ind$ will commute with inductive limits, we may assume that $\pi$ is actually a smooth representation of $H(O)_m=H(O)/H(O)(m+1)$. As before consider $\Gr_{H,m}=H(K)/H(O)(m+1)$. It has an action of $H(K)$ on the left and of $H(O)_m$ on the right.  We know that it can be written as an inductive limit of some $X_i$ which are finite-dimensional $F$-varieties where the maps $X_i\to X_{i+1}$ are closed embeddings. We may also assume that each $X_i$ is $H(O)_m$-invariant.

\begin{definition}\label{indHO}
We set $\mcS(X_i,V)$ to be the space of functions $f:X_i\to V$ satisfying the following conditions:

\begin{enumerate}
\item $f$ is locally constant
\item $f$ is $H(O)_m$-equivariant, i.e. $f(x h)=\pi(h)f(x)$ where $x\in X_i$ and $h\in H(O)_m$.
\item The quotient of the support of $f$ by $H(O)_m$ is compact.
\end{enumerate}
We have natural restriction maps $\mcS(X_{i+1},V)\to\mcS(X_i,V)$.
Thus we can form the pro-vector space $\mS(V)$ equal to the projective limit of the $\mcS(X_i,V)$.
\end{definition}
The pro-vector space $\mS(V)$ has an obvious action of the group $H(K)$ on the left (one has to check that this is actually a smooth representation of $H(K)$ on the pro-vector space $\mathbb S(V)$ -- this is done in Section 3 of \cite{GK}; also one has to check that it is independent of the choice of the $X_i$ which is easy and we shall not discuss it here).

The same definition makes sense for induction from any "thick" subgroup of $H(O)$ (cf. \cite{GK}): by definition a subgroup of $H(O)$ is thick if it is equal to the preimage of a closed subgroup of $H(O)_m$ for some $m\geq 0$.
\begin{remark}
The reader should note that this functor is denoted by $i\, _{H(O)}^{H(K)}$ in \cite{GK}. In {\em loc.cit.} the authors also define another induction functor $\widetilde i\, _{H(O)}^{H(K)}$ which is always right adjoint to the restriction functor from $H(K)$ to $H(O)$. This functor is endowed with a natural map of functors
$i\, _{H(O)}^{H(K)}\to \widetilde i\, _{H(O)}^{H(K)}$. The definition of $\widetilde i\, _{H(O)}^{H(K)}$ is similar to the above, except that we drop condition (3) above. In the case when $H$ is reductive, the quotient $X_i/H(O)_m$ is automatically compact, so this condition is satisfied automatically, therefore for reductive $H$ the two functors coincide. This is the reason why in this case our functor $\ind_{H(O)}^{H(K)}$ is actually right adjoint to the restriction functor.
\end{remark}

\subsection{The representations $\Pi(\pi)$} Let $\pi$ be a special representation of $G(O)$ of depth $m$. Then as above we can consider the induced representation $\Pi(\pi)=\ind_{G(O)}^{G(K)}\pi$. The main goal of this Section is to study the restriction of $\Pi(\pi)$ to $P(K)$ (for example we are going to show that this restriction is irreducible).

\subsection{A contradiction and its resolution} At this point it may seem that we have a contradiction. Namely, since
$G(K)=P(K)\cdot G(O)$, it looks as though Proposition \ref{indHL}, together with the description of induction from $P(O)$ given above, implies that the restriction of every $\Pi(\pi)$ to $P(K)$ would have to be isomorphic to $\Upsilon$. However, Corollary \ref{no-rest} says that this is impossible.

This contradiction is resolved as follows. The groups $G(K)$ and $P(K)$ are not algebraic $F$-groups (they are only group ind-schemes), so Proposition \ref{indHL} is not applicable in this situation. More precisely, as mentioned in Remark \ref{not-isom}, although the map $\Gr_P\to \Gr_G$ is bijective and continuous, it is not a homeomorphism: each $\Gr_P^i$ is closed in $\Gr_P$, whereas its image in $\Gr_G$, the semi-infinite orbit $S^i$, is not closed. Hence the analogue of Lemma \ref{open} fails here. This is a genuinely new feature of working with groups over $K$. Below we formulate a result that partially describes the restriction of $\Pi(\pi)$ to $P(K)$ and its relation to $\Upsilon$; see Remark \ref{ass-gr}.
%--------------------------------------------------------------------------

\subsection{Restricting $\Pi(\pi)$ to $P(K)$}
Let $\pi$ be as above. We denote by  $\Theta (\pi)$ the restriction of the representation $ \Pi (\pi)$ to $P(K)$. Below is the main result of this Section.
\begin{theorem}\label{exist}
\begin{enumerate}
\item For every special representation $\pi$ of $G(O)$, the representation $\Theta(\pi)$ is complete.
\item One has $\bar\Theta(\pi)\simeq\Upsilon$.
\end{enumerate}
\end{theorem}
\begin{claim}\label{no-Gplus-invariants}
For every special representation $\pi$ of $G(O)$, one has
$$
\Pi(\pi)^{G^+(K)}=0.
$$
The analogous statement holds for a special representation of $G'(O)$.
\end{claim}
\begin{corollary}\label{irrr}
The representation $\Theta(\pi)$ is irreducible. In particular, $\Pi(\pi)$ is irreducible.
\end{corollary}
\begin{proof}
Since $O\subset N(K)\subset G^+(K)$, Lemma \ref{Gplus-invariants} gives
$$
\Theta(\pi)^{N(K)}=\Theta(\pi)^O=\Pi(\pi)^{G^+(K)}.
$$
Thus Claim \ref{no-Gplus-invariants} implies that $\Theta(\pi)^{N(K)}=0$. Theorem \ref{exist}(1) shows that $\Theta(\pi)$ is complete, and Theorem \ref{exist}(2) gives $\bar\Theta(\pi)\simeq\Upsilon$. Proposition \ref{isbar} therefore implies that $\Theta(\pi)$ is irreducible. Since every $G(K)$-invariant subobject is, in particular, $P(K)$-invariant, $\Pi(\pi)$ is irreducible as well.
\end{proof}

\subsection{Proof of Theorem \ref{exist}}
We shall use the following consequence of the construction of coinvariants for quasi-pro-unipotent group objects.
\begin{lemma}\label{strict-limit-coinvariants}
Let $Q$ be quasi-pro-unipotent and let $\{V_i\}$ be a strict inverse system of smooth $Q$-representations with surjective transition maps. Then the natural morphism
$$
(\varprojlim_iV_i)_Q\longrightarrow\varprojlim_i(V_i)_Q
$$
is an isomorphism.
\end{lemma}
\begin{proof}
Write $Q=\varprojlim_rQ_r$, with $Q_r$ unipotent of finite type. At every finite stage, $Q_r$-coinvariants are exact and are computed by the cokernel of the action map. Since the transition maps of the system $\{V_i\}$ are surjective, the systems of kernels, images, and cokernels occurring in this presentation satisfy the Mittag--Leffler condition. Therefore taking $\varprojlim_i$ preserves the corresponding exact sequence and commutes with the cokernel. Passing to the projective limit over $r$ gives the asserted isomorphism. This is also the strict inverse-limit compatibility implicit in the construction of coinvariants for quasi-pro-unipotent group objects in \cite[Proposition~2.7]{GK1}.
\end{proof}
Let $\pi$ be any representation of $G(O)_m$, regarded as a representation of $G(O)$. Then the representation $\Pi(\pi)$ has a canonical increasing $\ZZ$-filtration
$$
\mcF^\sharp=\{\mF_i^\sharp\}_{i\in\ZZ}
$$
with the following properties:

1) $\mcF^\sharp$ is admissible.

2) Each $\mF_n^\sharp$ is $P_0$-invariant, and for every $r\in\ZZ$ one has
$$
\widehat t^{\,r}\mF_n^\sharp=\mF_{n-r}^\sharp.
$$
In particular, $P(K)$ acts on
$$
\bar\Theta(\pi)^\sharp:=\prod_{i\in\ZZ}\mF_{i+1}^\sharp/\mF_i^\sharp.
$$

3) $\bar\Theta(\pi)^\sharp$ is naturally isomorphic to $\ind_{P(O)}^{P(K)}(\pi|_{P(O)})$.

\begin{remark}\label{ass-gr}
As was mentioned above one might naively expect that for every $\pi$ we have $\Pi(\pi)|_{P(K)}=\ind_{P(O)}^{P(K)}(\pi|_{P(O)})$, but in practice it actually fails. Statement 3) above is the best substitute for this that we might hope for.
\end{remark}
The filtration $\mcF^\sharp$ is defined by setting $\mF_n^\sharp$ to be the projective limit of the spaces of functions from Definition \ref{indHO} (for $H=G$) supported on
$$
X_i\cap\left(\bigcup_{j\geq -n}S_m^j\right).
$$
This is an open subset of $X_i$. The filtration is increasing because the support condition becomes weaker as $n$ increases.
The properties 1 and 2 above are straightforward, property 3 follows immediately from the fact that $\Gr_{P,m}^j$ is a $P(O)_m$-torsor over $\Gr_P^j=S^j$ and the induced $G(O)_m$-torsor is isomorphic to $S^j_m$.

Thus to prove both assertions of Theorem \ref{exist} it is enough to show the following:
\begin{lemma}
Let $\pi$ be a special representation of $G(O)$ of depth $m$. Consider the filtrations $\mcF$ and $\mcF^{\sharp}$ on $\Pi(\pi)$.
Then
$$
\mF_n^\sharp=\mF_{m-n}.
$$
Equivalently, $\mF_{-n}^\sharp=\mF_{m+n}$.
\end{lemma}
\begin{proof}
Since the corresponding terms of both filtrations are permuted by the elements $\widehat t^{\,r}$, it is enough to prove
$$
\mF_m^\sharp=\mF_0.
$$
We may assume that the $X_i$ are $G(O)$-invariant. Then both $\mF_m^\sharp$ and $\mF_0$ are the projective limits of their images in $\mcS(X_i)$. It is therefore enough to prove that the kernel of
$$
\mcS(X_i)\longrightarrow\mcS(X_i)_O
$$
consists exactly of the functions supported on
$$
X_i\cap\left(\bigcup_{j\geq-m}S_m^j\right).
$$
Suppose that $f$ vanishes after taking $O$-coinvariants. Its restriction to
$$
X_i\cap\left(\bigcup_{j<-m}S_m^j\right)
$$
vanishes, since $O$ acts trivially on $\bigcup_{j<-m}S_m^j$. Hence $f$ is supported on $X_i\cap(\bigcup_{j\geq-m}S_m^j)$. This proves
$$
\mF_0\subset\mF_m^\sharp
$$
for every representation of $G(O)_m$, without any specialness assumption.

Now assume that $\pi$ is special of depth $m$. We need to prove that $(\mF_m^\sharp)_O=0$. It is enough to prove
$$
\left(\mF_{j+1}^\sharp/\mF_j^\sharp\right)_O=0
\qquad(j<m).
$$
Indeed, the functor of $O$-coinvariants is exact. For every $i$, the image of the filtration in $\mcS(X_i)$ has only finitely many nonzero graded pieces. Hence finite induction on these pieces shows that the $O$-coinvariants of the image of $\mF_m^\sharp$ in $\mcS(X_i)$ vanish. These images form a strict inverse system with surjective transition maps. Lemma \ref{strict-limit-coinvariants} therefore gives
$$
(\mF_m^\sharp)_O\simeq
\varprojlim_i(\operatorname{im}(\mF_m^\sharp\to\mcS(X_i)))_O=0.
$$

By property (3) above and Proposition \ref{Upsilonm-M}, the quotient
$$
\mF_{j+1}^\sharp/\mF_j^\sharp
$$
is isomorphic to $\mcM$ with the action $\rho_{m-j-1}$. Since $j<m$, one has $m-j-1\geq0$, and its $O$-coinvariants vanish. This proves $\mF_m^\sharp\subset\mF_0$, and therefore $\mF_m^\sharp=\mF_0$.

\end{proof}


%-----------------------------------------------------------------------------------------------
\subsection{Explicit realization of $\Pi(\pi)$} For completeness let us explain how to realize it in an explicit pro-vector space $\mW(\pi)$ for representations $\pi$ of $G(O)$ of odd depth. By Corollary~\ref{cunr}, write $\pi$ as the inflation of $\ind_R^{G(O)_m}(\nu)$, where $R\subset G(O)_m$ is a closed almost unipotent subgroup and $\nu$ is a character of $R$.

Consider the $F$-variety $\Gr_{G,m}^l$. It is acted on (on the right) by the group $G(O)_m$, hence by its subgroup $R$. Let us set
$$
W_l(\pi):= \mcS (\oGr_{G,m}^l)_{R,\nu}.
$$
It can also be described as the space of locally constant functions on $\oGr_{G,m}^l$ which are

a) Eigen-functions of $R$ with eigen-value $\nu$

b) Have compact support mod $R$.

\noindent
The embeddings $\oGr_{G,m}^l\hookrightarrow \oGr_{G,m}^{l+2}$ define restriction maps $W_{l+2}(\pi)\to W_l(\pi)$, and $\mW(\pi)$ is the pro-space equal to the projective limit of all the $W_l(\pi)$ (note that it is naturally a direct sum of two sub-pro-spaces corresponding to the parity of $l$, but this decomposition is not preserved by $G(K)$).

It is easy to see that
$$
\mW(\pi)=\mS(\Gr_{G,m})_{R,\nu},
$$
where $\mS(\Gr_{G,m})$ denotes the pro-vector space equal to the projective limit of all the $\mcS(\oGr_{G,m}^l)$.

\subsection{Representations of $G(K)$ induced from special representations of $G'(O)$}
We record the changes required for induction from $G'(O)$. They are important because the relevant depths are half-integral.

First of all we can define the induction functor $\ind_{G'(O)}^{G(K)}$ (again, this is a special case of the corresponding definition from \cite{GK} since $G'(O)$ is a thick subgroup in the terminology of {\em loc. cit.}). If $\pi$ is a special representation of $G'(O)$ we again set
$$
\Pi(\pi)=\ind_{G'(O)}^{G(K)}\pi.
$$
Let $m=n-\frac12$ be the depth of $\pi$. Use the lattice chain defining $G'(O)$ to index the twisted affine Grassmannian by $\frac12\ZZ$. Its semi-infinite strata are $(S^j_{n-1})'$, and conjugation by $\widehat t$ shifts their indices by one. Define the support filtration $\mcF^\sharp$ by the same support condition as above, with $S_m^j$ replaced throughout by $(S^j_{n-1})'$. The geometry of these strata gives analogues of properties (1)--(3): the filtration is admissible, it is shifted by $\widehat t$, and its associated graded is
$$
\ind_{P(O)}^{P(K)}(\pi|_{P(O)}).
$$
Since $\pi$ is special, its restriction to $P(O)$ is the special representation $\bs(n-1)$. Proposition \ref{Upsilonm-M} identifies each consecutive quotient below the critical stratum with $\mcM$ carrying an action $\rho_r$ with $r\geq0$. Its $O$-coinvariants therefore vanish. Exactness, finite induction over the strata at every finite-dimensional stage, and Lemma \ref{strict-limit-coinvariants} show that the support filtration agrees with the intrinsic coinvariant filtration, exactly as in the proof for $G(O)$. Thus the restriction $\Theta(\pi)$ is complete and
$$
\bar\Theta(\pi)\simeq\Upsilon.
$$
The proof that there are no $N(K)$-invariants is unchanged: $O\subset N(K)\subset G^+(K)$, the conjugates of $O$ generate $G^+(K)$, and the analogue of Claim \ref{no-Gplus-invariants} follows from the same twisted-stratum calculation. Proposition \ref{isbar} now shows that $\Theta(\pi)$ is irreducible. Hence $\Pi(\pi)$ is irreducible as well. This proves the analogue of Theorem \ref{exist} and Corollary \ref{irrr} for $G'(O)$.

\section{Questions}\label{quest}
Here we present a list of questions to which we do not know an answer and possible directions of future research.
 \begin{enumerate}

\item
Let $\pi, \pi'$ be two non-isomorphic special representations of either $G(O)$ or $G'(O)$ (i.e. either these are two representations of two different groups or two non-isomorphic representations of the same group).
Can it happen that $\Pi(\pi)$ and $\Pi(\pi')$ are isomorphic? 

\item
Is it true that any special representation of $G(K)$ is isomorphic to $\Pi(\pi)$ where $\pi$ is a special representation of either $G(O)$ or $G'(O)$?
\item
Let $\pi$ and $\pi'$ be two special representations of either $G(O)$ or $G'(O)$. Is it true that $\Theta(\pi)$ is isomorphic to $\Theta(\pi')$ if and only if $m(\pi)=m(\pi')$? (Note that since the depth of any special representation of $G(O)$ is an integer, and the depth of any special representation of $G'(O)$ is not an integer, the condition $m(\pi)=m(\pi')$ implies that these are representations of the same group).
\footnote{We expect that the answers to both Question 2 and Question 3 are positive, but we do not know how to prove this at the moment.}
\item
Does there exist a representation $\mW$ of $G(K)$ such that
$$
1<\dim\mcW(\mW)<\infty,
$$
where $\mcW(\mW):=\Hom_{N(K)}(\psi,\mW)$ is the Whittaker space of $\mW$?

\item
The definition of elliptic representations  extends naturally
to the case of maximal bounded  subgroups of  $G(K)$ for an arbitrary split reductive group $G$ (the definition is completely straightforward for the case of $G(O)$; for other maximal bounded subgroups it requires more work, which we shall not discuss here). Can one classify those? Is it true that their induction to $G(K)$ is irreducible?

\item
The notion of  special representations  extends verbatim to the case of  $\PGL(n)$ (in this case $P$ should be the maximal parabolic subgroup which is equal to the stabilizer of a point in $\mathbb P^{n-1}$). Is it possible to extend the results of this paper to the group $\PGL(n)$?

\item
Let $\mcH _{q,t}$ be the double affine Hecke algebra for
$\PGL(2)$. In Section 7.8 of \cite{GK1} there is a
construction of a pro-vector space  $\mW$ endowed with two commuting actions of a central extension of $G(K)$ and the algebra $\mcH_{q,t}$. Therefore we obtain
representations of the algebra $ \mcH _{q,t} $ on
$\mcW (\mW)$. How to describe this  representation?
\end{enumerate}

\begin{thebibliography}{99}
\bibitem{BZ}
J.~Bernstein and A.~Zelevinskii,
{\em Representations of the group $GL(n,F)$ where $F$ is a local non-Archimedean field},
Uspekhi Mat. Nauk 31:3 (1976), 5--70.
\bibitem{BD}
A.~Beilinson and V.~Drinfeld,
{\em Quantization of Hitchin's Hamiltonians and Hecke eigensheaves},
\url{https://math.uchicago.edu/~drinfeld/langlands/QuantizationHitchin.pdf}.
\bibitem{BT}
F.~Bruhat and J.~Tits,
{\em Groupes r\'eductifs sur un corps local},
Publ. Math. IH\'ES {\bf 41} (1972), 5--251.
\bibitem{BH}
C.~Bushnell and G.~Henniart,
{\em The Local Langlands Conjecture for $GL(2)$},
Grundlehren der mathematischen Wissenschaften, vol.~335, Springer-Verlag, Berlin, 2006.
\bibitem{GK}
D.~Gaitsgory and D.~Kazhdan,
{\em Representations of algebraic groups over a $2$-dimensional local field},
Geom. Funct. Anal. {\bf 14} (2004), no.~3, 535--574.
\bibitem{GK1}
D.~Gaitsgory and D.~Kazhdan,
{\em Algebraic groups over a $2$-dimensional local field: some further constructions},
in {\em Studies in Lie Theory}, Progress in Mathematics, vol.~243,
Birkh\"auser Boston, 2006, 97--130.
\bibitem{GK2}
I.~M.~Gelfand and D.~Kazhdan,
{\em Representations of the group $GL(n,K)$ where $K$ is a local field},
in {\em Lie Groups and Their Representations}, Halsted Press, 1975, 95--118.
\bibitem{HW}
T.~K.~Howard and M.~H.~Weissman,
{\em Depth-zero representations of nonlinear covers of $p$-adic groups},
Int. Math. Res. Not. IMRN 2009, no.~21, 3979--3995.
\bibitem{H}
B.~Hsie,
{\em Mackey theory for $p$-adic Lie groups},
Acta Math. Sin. (Engl. Ser.) {\bf 22} (2006), no.~2, 507--514.
\bibitem{M}
G.~W.~Mackey,
{\em Induced Representations of Groups and Quantum Mechanics},
W.~A. Benjamin, New York, 1968.
\bibitem{Moore}
C.~C.~Moore,
{\em Group extensions of $p$-adic and adelic linear groups},
Publ. Math. IH\'ES {\bf 35} (1968), 5--70.
\bibitem{MV}
I.~Mirkovi\'c and K.~Vilonen,
{\em Geometric Langlands duality and representations of algebraic groups over commutative rings},
Ann. of Math. (2) {\bf 166} (2007), no.~1, 95--143.
\bibitem{O}
R.~Olbricht,
{\em Local $\epsilon$-factors and the Langlands local correspondence}, preprint, 2005.
\bibitem{RangaRao}
R.~Ranga Rao,
{\em On some explicit formulas in the theory of Weil representation},
Pacific J. Math. {\bf 157} (1993), no.~2, 335--371.
\bibitem{Roberts}
B.~Roberts,
{\em The theta correspondence for similitudes},
Israel J. Math. {\bf 94} (1996), 285--317.
\bibitem{S}
J.-P.~Serre,
{\em Trees}, translated by J.~Stillwell,
Springer-Verlag, Berlin--New York, 1980.
\bibitem{T}
J.~Tits,
{\em Reductive groups over local fields},
in {\em Automorphic Forms, Representations and $L$-Functions},
Proc. Sympos. Pure Math., vol.~33, Part~1, Amer. Math. Soc., 1979, 29--69.
\bibitem{W}
A.~Weil,
{\em Sur certains groupes d'op\'erateurs unitaires},
Acta Math. {\bf 111} (1964), 143--211.
\bibitem{Yu}
J.-K.~Yu,
{\em Construction of tame supercuspidal representations},
J. Amer. Math. Soc. {\bf 14} (2001), no.~3, 579--622.
\end{thebibliography}
\end{document}
